%===========================================================================
\section*{Sources}
\phantomsection\addcontentsline{toc}{section}{Sources}
\label{sec:sources}
%===========================================================================

These are the works the book leans on, for their results, their history and the company they keep. Each one is listed under the first chapter that names it. A note after an entry says what the book takes from it, where the title doesn't already say. Where I couldn't confirm a detail of an entry, a page range or a place of printing, I left it out rather than guess.

\subsection*{Page one}

\noindent Alan M. Turing (1936--1937). On Computable Numbers, with an Application to the Entscheidungsproblem. Proceedings of the London Mathematical Society, series 2, volume 42, pp. 230--265. \href{https://doi.org/10.1112/plms/s2-42.1.230}{DOI: 10.1112/plms/s2-42.1.230}. The clerk at a desk reduced to a tape, a head, a state and a card of rules, with an appendix sketching the match with Church's version.
\par\medskip

\subsection*{Chapter 1. Any machine that steps}

\noindent American Institute of Physics. Harvard College Observatory: The labor and status of the women computers. \href{https://www.aip.org/history/harvard-observatory/labor-status-women-computers}{AIP page}. The women computers were paid a fraction of what the men were paid; the page puts the men's hourly rate at nearly five times as much.
\par\medskip

\noindent Alonzo Church (1936). An Unsolvable Problem of Elementary Number Theory. American Journal of Mathematics 58(2), pp. 345--363. \href{https://doi.org/10.2307/2371045}{DOI: 10.2307/2371045}. Computation written as functions handed to functions, from the same year as Turing's paper.
\par\medskip

\noindent B. Jack Copeland (1997, revised 2023). The Church-Turing Thesis. Stanford Encyclopedia of Philosophy. \href{https://plato.stanford.edu/entries/church-turing/}{Entry}. Why the claim about everything a person could do by following rules is a thesis and not a theorem.
\par\medskip

\noindent Martin Davis (1958). Computability and Unsolvability. New York: McGraw-Hill. Reprinted by Dover, 1982. The first use of the words ``the halting problem.''
\par\medskip

\noindent David Alan Grier (2005). When Computers Were Human. Princeton University Press. Rooms of human computers working out artillery tables in both world wars.
\par\medskip

\noindent Joel David Hamkins and Theodor Nenu (2026). Did Turing prove the undecidability of the halting problem? Journal of Logic and Computation 36(1), exaf075. \href{https://doi.org/10.1093/logcom/exaf075}{DOI: 10.1093/logcom/exaf075}; \href{https://arxiv.org/abs/2407.00680}{arXiv: 2407.00680}. What Turing proved in 1936, Kleene's statement of the halting problem in 1952, and Davis's name for it in 1958.
\par\medskip

\noindent Stephen Cole Kleene (1952). Introduction to Metamathematics. Amsterdam: North-Holland; New York: D. Van Nostrand. Where the halting problem is first stated in the form people now learn.
\par\medskip

\noindent Henrietta S. Leavitt and Edward C. Pickering (1912). Periods of 25 Variable Stars in the Small Magellanic Cloud. Harvard College Observatory Circular 173, pp. 1--3. The link between a variable star's period and its brightness that became a way to tell how far away a star is.
\par\medskip

\noindent J. C. Shepherdson and H. E. Sturgis (1963). Computability of Recursive Functions. Journal of the ACM 10(2), pp. 217--255. \href{https://doi.org/10.1145/321160.321170}{DOI: 10.1145/321160.321170}. Machines with registers instead of a tape compute the same things Turing's machine does.
\par\medskip

\noindent Dava Sobel (2016). The Glass Universe: How the Ladies of the Harvard Observatory Took the Measure of the Stars. Viking. The women who catalogued stars by hand at Harvard, Henrietta Leavitt among them.
\par\medskip

\noindent Alan M. Turing (1937). Computability and $\lambda$-definability. Journal of Symbolic Logic 2(4), pp. 153--163. \href{https://doi.org/10.2307/2268280}{DOI: 10.2307/2268280}. The full proof that Turing's machines and Church's functions compute the same things.
\par\medskip

\subsection*{Chapter 2. Running out of room}

\noindent Charles H. Bennett (1973). Logical reversibility of computation. IBM Journal of Research and Development 17(6), pp. 525--532. \href{https://doi.org/10.1147/rd.176.0525}{DOI: 10.1147/rd.176.0525}. A computation that saves its intermediate results can run without erasing and then undo its own steps.
\par\medskip

\noindent P. Bocchieri and A. Loinger (1957). Quantum Recurrence Theorem. Physical Review 107(2), pp. 337--338. \href{https://doi.org/10.1103/PhysRev.107.337}{DOI: 10.1103/PhysRev.107.337}. A closed quantum system comes back close to where it started.
\par\medskip

\noindent Ludwig Boltzmann (1896). Entgegnung auf die w\"armetheoretischen Betrachtungen des Hrn. E. Zermelo. Annalen der Physik und Chemie, Neue Folge 57, pp. 773--784. \href{https://doi.org/10.1002/andp.18962930414}{DOI: 10.1002/andp.18962930414}. The reply to Zermelo that nobody would ever wait long enough to see a return.
\par\medskip

\noindent Alex Brodsky and Nicholas Pippenger (2002). Characterizations of 1-way quantum finite automata. SIAM Journal on Computing 31(5), pp. 1456--1478. \href{https://doi.org/10.1137/S0097539799353443}{DOI: 10.1137/S0097539799353443}. A quantum automaton that measures once at the end recognises exactly the group languages; one that measures after every letter picks up ``has seen a b.''
\par\medskip

\noindent Richard Dedekind (1888). Was sind und was sollen die Zahlen? Braunschweig: Vieweg. The definition of infinite: a collection that matches one-to-one with a part of itself that leaves something out.
\par\medskip

\noindent David Deutsch (2013). Constructor theory. Synthese 190(18), pp. 4331--4359. \href{https://doi.org/10.1007/s11229-013-0279-z}{DOI: 10.1007/s11229-013-0279-z}. Physics stated as which transformations can and can't be brought about.
\par\medskip

\noindent David Deutsch and Chiara Marletto (2015). Constructor theory of information. Proceedings of the Royal Society A 471(2174), 20140540. \href{https://doi.org/10.1098/rspa.2014.0540}{DOI: 10.1098/rspa.2014.0540}; \href{https://arxiv.org/abs/1405.5563}{arXiv: 1405.5563}. Its section 8.6: measuring one quantity disturbs another, and nothing undoes the disturbance while the result is kept.
\par\medskip

\noindent Peter Gustav Lejeune Dirichlet (1842). Recherches sur les formes quadratiques \`a coefficients et \`a ind\'etermin\'ees complexes. Journal f\"ur die reine und angewandte Mathematik 24, pp. 291--371. \href{https://doi.org/10.1515/crll.1842.24.291}{DOI: 10.1515/crll.1842.24.291}. One of his two papers of 1842 that use the drawer principle.
\par\medskip

\noindent Peter Gustav Lejeune Dirichlet (1842). Verallgemeinerung eines Satzes aus der Lehre von den Kettenbr\"uchen nebst einigen Anwendungen auf die Theorie der Zahlen. Bericht \"uber die Verhandlungen der K\"oniglich Preussischen Akademie der Wissenschaften, pp. 93--95. The other 1842 paper that uses the drawer principle.
\par\medskip

\noindent Edward Fredkin and Tommaso Toffoli (1982). Conservative logic. International Journal of Theoretical Physics 21(3--4), pp. 219--253. \href{https://doi.org/10.1007/BF01857727}{DOI: 10.1007/BF01857727}. Circuits that compute anything and lose nothing, by carrying extra inputs in and extra outputs out.
\par\medskip

\noindent George Gamow (1947). One Two Three\ldots{} Infinity. New York: Viking Press. The book that made Hilbert's hotel famous.
\par\medskip

\noindent David Hilbert (2013). David Hilbert's Lectures on the Foundations of Arithmetic and Logic 1917--1933, edited by William Ewald and Wilfried Sieg. Berlin and Heidelberg: Springer. \href{https://doi.org/10.1007/978-3-540-69444-1}{DOI: 10.1007/978-3-540-69444-1}. Prints the notes of the winter 1924 to 1925 course on the infinite, where the hotel appears.
\par\medskip

\noindent Helge Kragh (2014). The True (?) Story of Hilbert's Infinite Hotel. \href{https://arxiv.org/abs/1403.0059}{arXiv: 1403.0059}. Takes back his earlier guess that Gamow invented the hotel, after Tilman Sauer pointed him to Hilbert's lecture notes.
\par\medskip

\noindent Rolf Landauer (1961). Irreversibility and Heat Generation in the Computing Process. IBM Journal of Research and Development 5(3), pp. 183--191. \href{https://doi.org/10.1147/rd.53.0183}{DOI: 10.1147/rd.53.0183}. Quoted: ``typically of the order of kT'' and ``show, or at least very strongly suggest,'' p.~183; ``no guarantee that this minimum is in fact achievable,'' p.~187; and on p.~187 the inputs that are all ONE, where ``there is no need to carry out any operation.'' \href{https://fab.cba.mit.edu/classes/862.22/notes/computation/Landauer-1961.pdf}{Paper scan}.
\par\medskip

\noindent Yves Lecerf (1963). Machines de Turing r\'eversibles. R\'ecursive insolubilit\'e en $n \in \mathbb{N}$ de l'\'equation $u = \theta^n u$, o\`u $\theta$ est un ``isomorphisme de codes.'' Comptes Rendus de l'Acad\'emie des Sciences de Paris 257, pp. 2597--2600. \href{https://revcomp.info/legacy/mpf/rc/Lecerf/english/english.html}{English translation}. Machines with no merging can imitate any Turing machine.
\par\medskip

\noindent Jean Leurechon (1622). Selectae Propositiones in Tota Sparsim Mathematica Pulcherrimae. Gaspard Bernard, p. 2. The one-line remark that two people must have the same number of hairs.
\par\medskip

\noindent Josef Loschmidt (1876). \"Uber den Zustand des W\"armegleichgewichtes eines Systems von K\"orpern mit R\"ucksicht auf die Schwerkraft. Sitzungsberichte der Kaiserlichen Akademie der Wissenschaften in Wien, Mathematisch-Naturwissenschaftliche Classe 73, Abt. II, from p. 128. The laws of motion run just as well backwards.
\par\medskip

\noindent Chiara Marletto, Vlatko Vedral, Laura T. Knoll, Fabrizio Piacentini, Ettore Bernardi, Enrico Rebufello, Alessio Avella, Marco Gramegna, Ivo Pietro Degiovanni and Marco Genovese (2022). Emergence of constructor-based irreversibility in quantum systems: theory and experiment. Physical Review Letters 128, 080401. \href{https://doi.org/10.1103/PhysRevLett.128.080401}{DOI: 10.1103/PhysRevLett.128.080401}; \href{https://arxiv.org/abs/2009.14649}{arXiv: 2009.14649}. An irreversible task built from quantum bits that partly swap with each other, run with photons.
\par\medskip

\noindent Kenichi Morita (1996). Universality of a reversible two-counter machine. Theoretical Computer Science 168(2), pp. 303--320. \href{https://doi.org/10.1016/S0304-3975(96)00081-3}{DOI: 10.1016/S0304-3975(96)00081-3}.
\par\medskip

\noindent Thomas E. Ouldridge and David H. Wolpert (2023). Thermodynamics of deterministic finite automata operating locally and periodically. New Journal of Physics 25, 123013. \href{https://doi.org/10.1088/1367-2630/ad1070}{DOI: 10.1088/1367-2630/ad1070}; \href{https://arxiv.org/abs/2208.06895}{arXiv: 2208.06895}. If the smallest machine for a job merges two states on some letter, every machine for it does, and running it gives off entropy; one worked example is the strings with no b.
\par\medskip

\noindent Jean-\'Eric Pin (1992). On reversible automata. In LATIN '92, 1st Latin American Symposium on Theoretical Informatics, Lecture Notes in Computer Science 583, Springer, pp. 401--416. \href{https://doi.org/10.1007/BFb0023844}{DOI: 10.1007/BFb0023844}. What automata whose letters never land two states on one can recognise.
\par\medskip

\noindent Henri Poincar\'e (1890). Sur le probl\`eme des trois corps et les \'equations de la dynamique. Acta Mathematica 13, pp. 1--270. The recurrence theorem.
\par\medskip

\noindent Beno\^it Rittaud and Albrecht Heeffer (2014). The Pigeonhole Principle, Two Centuries Before Dirichlet. The Mathematical Intelligencer 36(2), pp. 27--29. \href{https://doi.org/10.1007/s00283-013-9389-1}{DOI: 10.1007/s00283-013-9389-1}. Leurechon's remark of 1622, and Dirichlet's uses of the principle in 1842 and his later name for it.
\par\medskip

\noindent Morgan Rogers (2019). Toposes of Discrete Monoid Actions. \href{https://arxiv.org/abs/1905.10277}{arXiv: 1905.10277}.
\par\medskip

\noindent Morgan Rogers (2021). Toposes of Monoid Actions. \href{https://arxiv.org/abs/2112.10198}{arXiv: 2112.10198}.
\par\medskip

\noindent Tilman Sauer and Ulrich Majer, eds. (2009). David Hilbert's Lectures on the Foundations of Physics 1915--1927. Berlin and Heidelberg: Springer. \href{https://doi.org/10.1007/b12915}{DOI: 10.1007/b12915}. The volume of Hilbert's lectures that Sauer helped edit.
\par\medskip

\noindent Tommaso Toffoli (1980). Reversible computing. In Automata, Languages and Programming (ICALP 1980), Lecture Notes in Computer Science, Springer, pp. 632--644. \href{https://doi.org/10.1007/3-540-10003-2_104}{DOI: 10.1007/3-540-10003-2\_104}. The proof that reversible gates compute everything.
\par\medskip

\noindent Jos Uffink (2004, revised 2024). Boltzmann's Work in Statistical Physics. Stanford Encyclopedia of Philosophy. \href{https://plato.stanford.edu/entries/statphys-Boltzmann/}{Entry}. Zermelo as Planck's assistant in 1896, his careful logical argument, and Boltzmann's reply ``drenched in sarcasm and bitterness.''
\par\medskip

\noindent Maria Violaris and Chiara Marletto (2022). Irreversibility beyond Landauer's principle: is it possible to perform erasure using the quantum homogenizer? New Journal of Physics 24, 113030. \href{https://doi.org/10.1088/1367-2630/ac9e13}{DOI: 10.1088/1367-2630/ac9e13}; \href{https://arxiv.org/abs/2205.11310}{arXiv: 2205.11310}. A machine made of reversible pieces that, run in a cycle, turns a pure quantum bit mixed and never back, suggesting erasure may cost more than the usual bound.
\par\medskip

\noindent Ernst Zermelo (1896). Ueber einen Satz der Dynamik und die mechanische W\"armetheorie. Annalen der Physik und Chemie, Neue Folge 57, pp. 485--494. \href{https://doi.org/10.1002/andp.18962930314}{DOI: 10.1002/andp.18962930314}. Uses Poincar\'e's recurrence theorem to argue that nothing in such a system can climb forever.
\par\medskip

\noindent Ernst Zermelo (1908). Untersuchungen \"uber die Grundlagen der Mengenlehre I. Mathematische Annalen 65(2), pp. 261--281. \href{https://doi.org/10.1007/BF01449999}{DOI: 10.1007/BF01449999}. Set theory put on axioms.
\par\medskip

\subsection*{Chapter 3. Why charge for it}

\noindent Charles H. Bennett (2003). Notes on Landauer's principle, reversible computation, and Maxwell's Demon. Studies in History and Philosophy of Modern Physics 34(3), pp. 501--510. \href{https://doi.org/10.1016/S1355-2198(03)00039-X}{DOI: 10.1016/S1355-2198(03)00039-X}; \href{https://arxiv.org/abs/physics/0210005v2}{arXiv: physics/0210005v2}. Quoted from the arXiv preprint: Earman and Norton's objection as ``in my opinion the one of greatest merit'' (p.~1); ballistic machines that compute ``at constant velocity without dissipating any energy'' and ``must be isolated'' from heat baths and noise, and the thermodynamics of error correction as ``a subject that appears to be in its infancy.''
\par\medskip

\noindent Ludwig Boltzmann (1877). \"Uber die Beziehung zwischen dem zweiten Hauptsatze der mechanischen W\"armetheorie und der Wahrscheinlichkeitsrechnung respektive den S\"atzen \"uber das W\"armegleichgewicht. Sitzungsberichte der Kaiserlichen Akademie der Wissenschaften in Wien, Mathematisch-Naturwissenschaftliche Classe 76, pp. 373--435. Entropy as the logarithm of the number of ways the parts could be arranged.
\par\medskip

\noindent Philippe Faist, Fr\'ed\'eric Dupuis, Jonathan Oppenheim and Renato Renner (2015). The minimal work cost of information processing. Nature Communications 6, 7669. \href{https://doi.org/10.1038/ncomms8669}{DOI: 10.1038/ncomms8669}; \href{https://arxiv.org/abs/1211.1037}{arXiv: 1211.1037}. The least work a logical step needs when it has to succeed every time, with an AND gate at $\log_2 3 \approx 1.6$.
\par\medskip

\noindent G. H. Hardy (1940). A Mathematician's Apology. Cambridge University Press. Mathematical reality lies outside the people who study it, and their theorems are notes on what they saw.
\par\medskip

\noindent Max Planck (1901). Ueber das Gesetz der Energieverteilung im Normalspectrum. Annalen der Physik 309(3), pp. 553--563. \href{https://doi.org/10.1002/andp.19013090310}{DOI: 10.1002/andp.19013090310}. Where the constant k and the entropy written as k times a logarithm come from.
\par\medskip

\noindent Max Planck (1920). The genesis and present state of development of the quantum theory. Nobel Lecture, 2 June 1920. In the translation by H. T. Clarke and L. Silberstein, The Origin and Development of the Quantum Theory, Oxford: Clarendon Press, 1922. \href{https://www.gutenberg.org/cache/epub/33663/pg33663-images.html}{1922 translation}. The constant ``is frequently termed Boltzmann's constant, although to the best of my knowledge Boltzmann himself never introduced it.''
\par\medskip

\noindent Heinrich Weber (1893). Leopold Kronecker. Mathematische Annalen 43, pp. 1--25. \href{https://doi.org/10.1007/BF01446613}{DOI: 10.1007/BF01446613}. Also printed in the Jahresbericht der Deutschen Mathematiker-Vereinigung 2. The memorial piece that carries the line about God and the whole numbers, pointing back to a lecture of 1886.
\par\medskip

\subsection*{Chapter 4. The smallest one}

\noindent Michael T. Clanchy (1979). From Memory to Written Record: England 1066--1307. London: Edward Arnold; third edition, Wiley-Blackwell, 2012. England learning to trust writing, and the forged seals and charters along the way.
\par\medskip

\noindent Stephen A. Cook (1971). The complexity of theorem-proving procedures. Proceedings of the Third Annual ACM Symposium on Theory of Computing (STOC '71), pp. 151--158. \href{https://doi.org/10.1145/800157.805047}{DOI: 10.1145/800157.805047}. A claim that comes with a certificate that's quick to check.
\par\medskip

\noindent Nicholas of Cusa (1433). De concordantia catholica. Where he calls the Donation of Constantine a forgery.
\par\medskip

\noindent Denise Schmandt-Besserat. Tokens and Writing: the Cognitive Development. Author's page, University of Texas at Austin. \href{https://sites.utexas.edu/dsb/tokens/tokens-and-writing-the-cognitive-development/}{Web page}. Hollow clay balls holding tokens, with the tokens pressed into the outside, around 3300 BC.
\par\medskip

\noindent Deuteronomy 19:15. \href{https://www.biblegateway.com/passage/?search=Deuteronomy+19:15&version=NRSV}{NRSV text}. Nobody is convicted on one witness; it takes two or three.
\par\medskip

\noindent Cynthia Dwork and Moni Naor (1992). Pricing via Processing or Combatting Junk Mail. Advances in Cryptology, CRYPTO '92, Lecture Notes in Computer Science 740 (Springer, 1993), pp. 139--147. \href{https://doi.org/10.1007/3-540-48071-4_10}{DOI: 10.1007/3-540-48071-4\_10}. Proof of work, proposed to make junk mail expensive.
\par\medskip

\noindent Robert W. Floyd (1967). Assigning meanings to programs. In J. T. Schwartz (ed.), Mathematical Aspects of Computer Science, Proceedings of Symposia in Applied Mathematics 19, American Mathematical Society, pp. 19--32. \href{https://doi.org/10.1090/psapm/019/0235771}{DOI: 10.1090/psapm/019/0235771}. A program stops because no infinite decreasing sequence exists in a well-ordered set.
\par\medskip

\noindent Forrester (n.d.). Books as Evidence, before and after Pacioli (1494 AD). Paper from the University of Strathclyde accounting history pages; the page gives no initials or date. \href{http://www.cs.trinity.edu/rjensen/readings/history/forrester/a01before1494AD.htm}{Mirror of the paper}. Reads Pacioli as reporting that books were taken to a consul to be authenticated before use.
\par\medskip

\noindent Dean P. Foster and Robert A. Stine (2008). $\alpha$-investing: a procedure for sequential control of expected false discoveries. Journal of the Royal Statistical Society, Series B 70(2), pp. 429--444. \href{https://doi.org/10.1111/j.1467-9868.2007.00643.x}{DOI: 10.1111/j.1467-9868.2007.00643.x}. A budget of false alarms that every test spends and every discovery earns back.
\par\medskip

\noindent Fox-IT (2012). Black Tulip: Report of the investigation into the DigiNotar Certificate Authority breach. Delft, 13 August 2012, for the Dutch Ministry of the Interior and Kingdom Relations. \href{https://zoek.officielebekendmakingen.nl/blg-191225.pdf}{Report}. Intruders took over the servers that signed certificates, and the company went bankrupt.
\par\medskip

\noindent Stuart Haber and W. Scott Stornetta (1991). How to Time-Stamp a Digital Document. Journal of Cryptology 3(2), pp. 99--111. \href{https://doi.org/10.1007/BF00196791}{DOI: 10.1007/BF00196791}. Records chained so that changing an old one breaks every one after it.
\par\medskip

\noindent C. A. R. Hoare (1969). An Axiomatic Basis for Computer Programming. Communications of the ACM 12(10), pp. 576--580. \href{https://doi.org/10.1145/363235.363259}{DOI: 10.1145/363235.363259}.
\par\medskip

\noindent Horace (c. 20 BC). Epistles, book 1, epistle 1, line 14. \href{https://www.thelatinlibrary.com/horace/epist1.shtml}{Latin text}. ``Nullius addictus iurare in verba magistri,'' the line the Royal Society cut its motto from.
\par\medskip

\noindent Gaius. Institutes, book 2, sections 104 and 119 (second century). \href{https://droitromain.univ-grenoble-alpes.fr/Anglica/gai2_Poste.htm}{Text and translation}. The will by bronze and balance before five witnesses (104), and a will sealed by seven witnesses (119).
\par\medskip

\noindent U.S. Department of Justice (2002). Statement of Deputy Attorney General Larry Thompson on the Arthur Andersen verdict, 15 June 2002. \href{https://www.justice.gov/archive/opa/pr/2002/June/02_dag_356.htm}{Press release}. Andersen convicted of obstruction of justice over the Enron documents.
\par\medskip

\noindent U.S. Department of Justice (2002). Press release, October 2002, on the criminal complaint against Andrew Fastow. \href{https://www.justice.gov/archive/opa/pr/2002/October/02_crm_568.htm}{Press release}. Enron's special purpose entities, set up to keep debt off the balance sheet.
\par\medskip

\noindent Geoffrey A. Lee (1977). The Coming of Age of Double Entry: The Giovanni Farolfi Ledger of 1299--1300. Accounting Historians Journal 4(2), pp. 79--96. \href{https://doi.org/10.2308/0148-4184.4.2.79}{DOI: 10.2308/0148-4184.4.2.79}. The Provence branch ledger of a Florentine company, the earliest complete double entry found.
\par\medskip

\noindent Henry Sumner Maine (1861). Ancient Law. London: John Murray, chapter VI. \href{https://www.gutenberg.org/ebooks/22910}{Project Gutenberg 22910}. The seven seals of the Roman will as the seven people of the older mock sale.
\par\medskip

\noindent Alvaro Martinelli (1983). The Ledger of Cristianus Lomellinus and Dominicus De Garibaldo, Stewards of the City of Genoa (1340--41). Abacus 19(2), pp. 83--118. \href{https://doi.org/10.1111/j.1467-6281.1983.tb00242.x}{DOI: 10.1111/j.1467-6281.1983.tb00242.x}. The Genoa treasurers' ledger for 1340, kept in double entry.
\par\medskip

\noindent Ralph C. Merkle (1979). Secrecy, Authentication, and Public Key Systems. PhD thesis, Stanford University. \href{https://www.ralphmerkle.com/papers/Thesis1979.pdf}{Thesis}. Where tree authentication, now called the Merkle tree, is introduced.
\par\medskip

\noindent Satoshi Nakamoto (2008). Bitcoin: A Peer-to-Peer Electronic Cash System. \href{https://bitcoin.org/bitcoin.pdf}{Paper}. Three of its eight references are to Haber and Stornetta.
\par\medskip

\noindent Luca Pacioli (1494). Summa de arithmetica, geometria, proportioni et proportionalit\`a. Venice. The bookkeeping treatise is the part titled Particularis de computis et scripturis.
\par\medskip

\noindent Peter E. Prosser (2001). Church History's Biggest Hoax. Christian History 72. \href{https://christianhistoryinstitute.org/magazine/article/church-historys-biggest-hoax}{Article}. The Donation most likely written in the eighth century, Cusa's exposure in 1433, and Valla's in 1440.
\par\medskip

\noindent Denise Schmandt-Besserat (1992). Before Writing, volume 1: From Counting to Cuneiform. Austin: University of Texas Press.
\par\medskip

\noindent Securities and Exchange Commission (1997). Reporting Requirements for Brokers or Dealers Under the Securities Exchange Act of 1934. Release No. 34-38245, Federal Register 62, pp. 6469--6474, 12 February 1997. \href{https://www.govinfo.gov/content/pkg/FR-1997-02-12/html/97-3426.htm}{Federal Register text}. Broker-dealer records kept electronically had to be on non-rewriteable, non-erasable media.
\par\medskip

\noindent Securities and Exchange Commission (2022). Electronic Recordkeeping Requirements for Broker-Dealers, Security-Based Swap Dealers, and Major Security-Based Swap Participants. Release No. 34-96034. \href{https://www.sec.gov/rules/final/2022/34-96034.pdf}{Final rule}. Adds an audit-trail alternative to write once, read many.
\par\medskip

\noindent Steven Shapin (1984). Pump and Circumstance: Robert Boyle's Literary Technology. Social Studies of Science 14(4), pp. 481--520. \href{https://doi.org/10.1177/030631284014004001}{DOI: 10.1177/030631284014004001}. Virtual witnessing, and the air pump as seventeenth-century big science.
\par\medskip

\noindent Supreme Court of the United States (2005). Arthur Andersen LLP v. United States, 544 U.S. 696. The conviction reversed over the jury instructions.
\par\medskip

\noindent The Royal Society. History of the Royal Society. \href{https://royalsociety.org/about-us/who-we-are/history/}{Web page}. First meeting 28 November 1660, and the motto Nullius in verba.
\par\medskip

\noindent Lorenzo Valla (c. 1440). De falso credita et ementita Constantini donatione declamatio. English translation: The Treatise of Lorenzo Valla on the Donation of Constantine, trans. Christopher B. Coleman, Yale University Press, 1922. \href{https://history.hanover.edu/texts/vallaintro.html}{Coleman translation}. The Donation taken apart by its Latin, ``satraps'' among the words no Roman of Constantine's day would have used.
\par\medskip

\subsection*{Chapter 5. More settled}

\noindent Lindsey Kuper and Ryan R. Newton (2013). LVars: lattice-based data structures for deterministic parallelism. Proceedings of the 2nd ACM SIGPLAN Workshop on Functional High-Performance Computing (FHPC 2013), pp. 71--84. \href{https://doi.org/10.1145/2502323.2502326}{DOI: 10.1145/2502323.2502326}. Memory cells whose value only climbs, where a program can only ask whether it has passed a point.
\par\medskip

\noindent Amin Timany and Lars Birkedal (2021). Reasoning about monotonicity in separation logic. Proceedings of the 10th ACM SIGPLAN International Conference on Certified Programs and Proofs (CPP 2021), pp. 91--104. \href{https://doi.org/10.1145/3437992.3439931}{DOI: 10.1145/3437992.3439931}. Any order carried in a program logic as a value that only climbs and that the program never reads.
\par\medskip

\subsection*{Chapter 6. Looking isn't keeping}

\noindent Charles H. Bennett (1982). The thermodynamics of computation: a review. International Journal of Theoretical Physics 21(12), pp. 905--940. \href{https://doi.org/10.1007/BF02084158}{DOI: 10.1007/BF02084158}. The demon taken apart one step at a time, with the cost moved from the look to the wipe.
\par\medskip

\noindent L\'eon Brillouin (1951). Maxwell's Demon Cannot Operate: Information and Entropy. I. Journal of Applied Physics 22(3), pp. 334--337. \href{https://doi.org/10.1063/1.1699951}{DOI: 10.1063/1.1699951}. To see a molecule the demon has to bounce light off it, and that costs entropy.
\par\medskip

\noindent Jeffrey Bub (2001). Maxwell's Demon and the Thermodynamics of Computation. Studies in History and Philosophy of Modern Physics 32(4), pp. 569--579. \href{https://doi.org/10.1016/S1355-2198(01)00023-5}{DOI: 10.1016/S1355-2198(01)00023-5}. The summary of Szilard and Brillouin: finding out costs at least $k \log 2$ of entropy for every bit.
\par\medskip

\noindent Arthur S. Eddington (1928). The Nature of the Physical World. Cambridge University Press. Chapter IV, The Running-Down of the Universe, where a theory against the second law has nothing for it but to ``collapse in deepest humiliation.''
\par\medskip

\noindent Paul N. Fahn (1996). Maxwell's demon and the entropy cost of information. Foundations of Physics 26(1), pp. 71--93. \href{https://doi.org/10.1007/BF02058888}{DOI: 10.1007/BF02058888}. Whether the look or the wipe gets charged is partly a choice of bookkeeping, and the total is what's fixed.
\par\medskip

\noindent Aur\'elien Garivier, Tor Lattimore and Emilie Kaufmann (2016). On explore-then-commit strategies. Advances in Neural Information Processing Systems 29 (NeurIPS 2016), pp. 784--792. \href{https://arxiv.org/abs/1605.08988}{arXiv: 1605.08988}. Over a long run, the best strategies that commit lose about twice as much as the best ones that keep adjusting.
\par\medskip

\noindent Sebastian Goldt and Udo Seifert (2017). Stochastic thermodynamics of learning. Physical Review Letters 118(1), 010601. \href{https://doi.org/10.1103/PhysRevLett.118.010601}{DOI: 10.1103/PhysRevLett.118.010601}. The entropy a learning network produces limits how much it can learn.
\par\medskip

\noindent Edwin T. Jaynes (1965). Gibbs vs Boltzmann entropies. American Journal of Physics 33(5), pp. 391--398. \href{https://doi.org/10.1119/1.1971557}{DOI: 10.1119/1.1971557}. Entropy as ``an anthropomorphic concept,'' a phrase he credits to Wigner (p.~398).
\par\medskip

\noindent Edwin T. Jaynes (1989). Clearing up Mysteries: The Original Goal. In John Skilling (ed.), Maximum Entropy and Bayesian Methods, Kluwer, pp. 1--27. \href{https://doi.org/10.1007/978-94-015-7860-8_1}{DOI: 10.1007/978-94-015-7860-8\_1}. Where the mind projection fallacy gets its name.
\par\medskip

\noindent Ruth E. Kastner (2025). Maxwell's Demon Is Foiled by the Entropy Cost of Measurement, Not Erasure. Foundations 5(2), article 16. \href{https://doi.org/10.3390/foundations5020016}{DOI: 10.3390/foundations5020016}; \href{https://arxiv.org/abs/2503.18186}{arXiv: 2503.18186}. Brillouin was essentially right, and it's the measuring that pays.
\par\medskip

\noindent James Clerk Maxwell (1865). A Dynamical Theory of the Electromagnetic Field. Philosophical Transactions of the Royal Society of London 155, pp. 459--512. \href{https://doi.org/10.1098/rstl.1865.0008}{DOI: 10.1098/rstl.1865.0008}.
\par\medskip

\noindent James Clerk Maxwell (1867, 1871). Letter to Peter Guthrie Tait, 11 December 1867, printed in Cargill Gilston Knott, Life and Scientific Work of Peter Guthrie Tait, Cambridge University Press, 1911, pp. 213--214 (``very observant and neat-fingered being,'' p.~214), with the undated note ``Concerning Demons'' on pp. 214--215 (``Call him no more a demon but a valve,'' p.~215); and Theory of Heat, Longmans, Green, London, 1871, the section on the limitation of the second law of thermodynamics.
\par\medskip

\noindent John D. Norton (2013). The end of the thermodynamics of computation: a no-go result. Philosophy of Science 80(5), pp. 1182--1192. \href{https://doi.org/10.1086/673714}{DOI: 10.1086/673714}. At the scale of molecules, thermal jiggling keeps a supposedly reversible process from finishing, and finishing it costs more than Landauer's amount.
\par\medskip

\noindent Thomas E. Ouldridge and Pieter Rein ten Wolde (2017). Fundamental Costs in the Production and Destruction of Persistent Polymer Copies. Physical Review Letters 118(15), 158103. \href{https://doi.org/10.1103/PhysRevLett.118.158103}{DOI: 10.1103/PhysRevLett.118.158103}.
\par\medskip

\noindent Thomas E. Ouldridge, Christopher C. Govern and Pieter Rein ten Wolde (2017). Thermodynamics of Computational Copying in Biochemical Systems. Physical Review X 7(2), 021004. \href{https://doi.org/10.1103/PhysRevX.7.021004}{DOI: 10.1103/PhysRevX.7.021004}. The same kind of accounting for how a cell reads out its own sensors.
\par\medskip

\noindent Jenny M. Poulton, Pieter Rein ten Wolde and Thomas E. Ouldridge (2019). Nonequilibrium correlations in minimal dynamical models of polymer copying. Proceedings of the National Academy of Sciences 116(6), pp. 1946--1951. \href{https://doi.org/10.1073/pnas.1808775116}{DOI: 10.1073/pnas.1808775116}; \href{https://arxiv.org/abs/1805.08502}{arXiv: 1805.08502}. When a copy grows with the weakest push, its advantage over random copying and the efficiency of storing it both vanish.
\par\medskip

\noindent Hans Reichenbach (1956). The Direction of Time. University of California Press, Berkeley. Rovelli credits him as the first to notice that any measurement dissipates.
\par\medskip

\noindent Richard Rhodes (1986). The Making of the Atomic Bomb. Simon and Schuster, New York. Chapter 1, ``Moonshine,'' opens with Szilard's own later account of the traffic light on Southampton Row and Rutherford's dismissal in the newspaper.
\par\medskip

\noindent L\'idia del Rio, Johan \AA{}berg, Renato Renner, Oscar Dahlsten and Vlatko Vedral (2011). The thermodynamic meaning of negative entropy. Nature 474(7349), pp. 61--63. \href{https://doi.org/10.1038/nature10123}{DOI: 10.1038/nature10123}. The cost of erasing depends on what the eraser already knows, down to zero, and below zero in the quantum case.
\par\medskip

\noindent Carlo Rovelli (2024). The Thermodynamic Cost of Choosing. Foundations of Physics 54(3), article 28. \href{https://doi.org/10.1007/s10701-024-00765-7}{DOI: 10.1007/s10701-024-00765-7}; \href{https://arxiv.org/abs/2309.08557}{arXiv: 2309.08557}. Quoted from the arXiv preprint: ``stabilization must be irreversible, otherwise it is not stabilization,'' with the thermostat set aside as predictable and ``Wahl ist Leiden'' near the end.
\par\medskip

\noindent Takahiro Sagawa and Masahito Ueda (2009). Minimal Energy Cost for Thermodynamic Information Processing: Measurement and Information Erasure. Physical Review Letters 102(25), 250602. \href{https://doi.org/10.1103/PhysRevLett.102.250602}{DOI: 10.1103/PhysRevLett.102.250602}; \href{https://arxiv.org/abs/0809.4098v3}{arXiv: 0809.4098v3}. A bound on the total cost of measuring and erasing, with the split depending on the memory; see also the erratum, Physical Review Letters 106, 189901 (2011).
\par\medskip

\noindent American Physical Society (2006). May, 1911: Rutherford and the Discovery of the Atomic Nucleus. APS News 15(5), This Month in Physics History. \href{https://www.aps.org/publications/apsnews/200605/history.cfm}{Article}.
\par\medskip

\noindent Leo Szilard (1929). \"Uber die Entropieverminderung in einem thermodynamischen System bei Eingriffen intelligenter Wesen. Zeitschrift f\"ur Physik 53(11--12), pp. 840--856. \href{https://doi.org/10.1007/BF01341281}{DOI: 10.1007/BF01341281}. The one-molecule engine.
\par\medskip

\noindent William Thomson (1874). Kinetic Theory of the Dissipation of Energy. Nature 9(232), pp. 441--444. \href{https://doi.org/10.1038/009441c0}{DOI: 10.1038/009441c0}. Where Maxwell's being first gets called a demon.
\par\medskip

\noindent Shoichi Toyabe, Takahiro Sagawa, Masahito Ueda, Eiro Muneyuki and Masaki Sano (2010). Experimental demonstration of information-to-energy conversion and validation of the generalized Jarzynski equality. Nature Physics 6(12), pp. 988--992. \href{https://doi.org/10.1038/nphys1821}{DOI: 10.1038/nphys1821}. A particle walked up a staircase by feedback on what a camera saw.
\par\medskip

\subsection*{Chapter 7. Looking through a window}

\noindent Edwin A. Abbott (1884). Flatland: A Romance of Many Dimensions, by A Square. Seeley and Co., London.
\par\medskip

\noindent MacTutor History of Mathematics. Edwin Abbott Abbott (J. J. O'Connor and E. F. Robertson). University of St Andrews. \href{https://mathshistory.st-andrews.ac.uk/Biographies/Abbott/}{Biography}. His mother was his father's first cousin, he was headmaster of the City of London School from 1865 to 1889, and his Times obituary didn't mention Flatland.
\par\medskip

\noindent Robert J. Aumann (1976). Agreeing to disagree. The Annals of Statistics 4(6), pp. 1236--1239. \href{https://doi.org/10.1214/aos/1176343654}{DOI: 10.1214/aos/1176343654}. Knowledge modelled by piles of states an agent can't tell apart.
\par\medskip

\noindent Fran\c{c}ois Bancilhon and Nicolas Spyratos (1981). Update semantics of relational views. ACM Transactions on Database Systems 6(4), pp. 557--575. \href{https://doi.org/10.1145/319628.319634}{DOI: 10.1145/319628.319634}.
\par\medskip

\noindent Mihir Bellare and Bennet S. Yee (1997). Forward integrity for secure audit logs. Technical report, Department of Computer Science and Engineering, University of California, San Diego. Where forward integrity gets its name.
\par\medskip

\noindent Laurent Bindschaedler, Quentin Botha and Christoph Siebenbrunner (2026). Agent Flight Recorder: Tamper-Evident Audit Trails with On-Chain Anchoring for Long-Horizon Tool-Using Agents. Preprint. \href{https://arxiv.org/abs/2609.01931}{arXiv: 2609.01931}.
\par\medskip

\noindent Manuel Blum (1983). Coin Flipping by Telephone: A Protocol for Solving Impossible Problems. ACM SIGACT News 15(1), pp. 23--27. \href{https://doi.org/10.1145/1008908.1008911}{DOI: 10.1145/1008908.1008911}. Presented at CRYPTO 1981; one side commits first and opens after the other calls, with the binding resting on how hard it is to factor.
\par\medskip

\noindent Jesse Bull and Joel Watson (2007). Hard evidence and mechanism design. Games and Economic Behavior 58(1), pp. 75--93. \href{https://doi.org/10.1016/j.geb.2006.03.003}{DOI: 10.1016/j.geb.2006.03.003}.
\par\medskip

\noindent Tyler Burge (1993). Content preservation. The Philosophical Review 102(4), pp. 457--488. \href{https://doi.org/10.2307/2185680}{DOI: 10.2307/2185680}. The acceptance principle: you're entitled to accept what you're told unless you have a reason not to.
\par\medskip

\noindent Michael Burrows, Mart\'in Abadi and Roger Needham (1990). A logic of authentication. ACM Transactions on Computer Systems 8(1), pp. 18--36. \href{https://doi.org/10.1145/77648.77649}{DOI: 10.1145/77648.77649}. If I believe an authority has the say over a statement and believe it holds the statement, I believe the statement.
\par\medskip

\noindent Alessandro Chiesa and Eran Tromer (2010). Proof-carrying data and hearsay arguments from signature cards. Proceedings of Innovations in Computer Science (ICS 2010), Tsinghua University Press, Beijing, pp. 310--331. \href{https://conference.iiis.tsinghua.edu.cn/ICS2010/content/papers/25.html}{Proceedings page}. Proofs passed along a chain of parties who don't trust each other.
\par\medskip

\noindent Robin Cockett, Isabelle Jianing Geng, Carlo Maria Scandolo and Priyaa Varshinee Srinivasan (2023). Extending resource monotones using Kan extensions. Proceedings of the Fifth International Conference on Applied Category Theory (ACT 2022), Electronic Proceedings in Theoretical Computer Science 380, pp. 203--223. \href{https://doi.org/10.4204/EPTCS.380.12}{DOI: 10.4204/EPTCS.380.12}. ``Some member'' and ``every member'' used to stretch measures of resources along a map.
\par\medskip

\noindent Edgar F. Codd (1970). A Relational Model of Data for Large Shared Data Banks. Communications of the ACM 13(6), pp. 377--387. \href{https://doi.org/10.1145/362384.362685}{DOI: 10.1145/362384.362685}.
\par\medskip

\noindent Patrick Cousot and Radhia Cousot (1977). Abstract interpretation: a unified lattice model for static analysis of programs by construction or approximation of fixpoints. Proceedings of the 4th ACM Symposium on Principles of Programming Languages (POPL '77), pp. 238--252. \href{https://doi.org/10.1145/512950.512973}{DOI: 10.1145/512950.512973}.
\par\medskip

\noindent Claude Cr\'epeau, Rafael Dowsley and Anderson C. A. Nascimento (2020). On the commitment capacity of unfair noisy channels. IEEE Transactions on Information Theory 66(6), pp. 3745--3752. \href{https://doi.org/10.1109/TIT.2020.2967048}{DOI: 10.1109/TIT.2020.2967048}; \href{https://arxiv.org/abs/1905.10921}{arXiv: 1905.10921}. How the rate drops when the one committing can steer the noise.
\par\medskip

\noindent Kenny Easwaran (2009). Probabilistic proofs and transferability. Philosophia Mathematica (III) 17(3), pp. 341--362. \href{https://doi.org/10.1093/philmat/nkn032}{DOI: 10.1093/philmat/nkn032}. A step-by-step proof carries its reason across to the reader, and a proof based on chance doesn't.
\par\medskip

\noindent Ronald Fagin, Joseph Y. Halpern, Yoram Moses and Moshe Y. Vardi (1995). Reasoning About Knowledge. MIT Press, Cambridge, Massachusetts.
\par\medskip

\noindent Joseph A. Goguen and Jos\'e Meseguer (1982). Security policies and security models. Proceedings of the 1982 IEEE Symposium on Security and Privacy, pp. 11--20. \href{https://doi.org/10.1109/SP.1982.10014}{DOI: 10.1109/SP.1982.10014}. Noninterference.
\par\medskip

\noindent Sanford C. Goldberg (2010). Relying on Others: An Essay in Epistemology. Oxford University Press. The weight put on the process behind the speaker's word.
\par\medskip

\noindent Shafi Goldwasser, Silvio Micali and Charles Rackoff (1985). The Knowledge Complexity of Interactive Proof-Systems. Proceedings of the 17th Annual ACM Symposium on Theory of Computing (STOC '85), pp. 291--304. \href{https://doi.org/10.1145/22145.22178}{DOI: 10.1145/22145.22178}.
\par\medskip

\noindent Jerry R. Green and Jean-Jacques Laffont (1986). Partially verifiable information and mechanism design. Review of Economic Studies 53(3), pp. 447--456. \href{https://doi.org/10.2307/2297639}{DOI: 10.2307/2297639}.
\par\medskip

\noindent Andreas Haeberlen, Petr Kouznetsov and Peter Druschel (2007). PeerReview: practical accountability for distributed systems. Proceedings of the 21st ACM Symposium on Operating Systems Principles (SOSP '07), pp. 175--188. \href{https://doi.org/10.1145/1294261.1294279}{DOI: 10.1145/1294261.1294279}. Every node keeps a secure log, and a checker replays it against a copy of the rules.
\par\medskip

\noindent Andreas Haeberlen and Petr Kuznetsov (2009). The fault detection problem. Principles of Distributed Systems, OPODIS 2009, Lecture Notes in Computer Science 5923, Springer, pp. 99--114. \href{https://doi.org/10.1007/978-3-642-10877-8_10}{DOI: 10.1007/978-3-642-10877-8\_10}. Which kinds of cheating an observer holding only the messages can catch at all.
\par\medskip

\noindent Joseph Y. Halpern, Yoram Moses and Mark R. Tuttle (1988). A knowledge-based analysis of zero knowledge. Proceedings of the 20th Annual ACM Symposium on Theory of Computing (STOC '88), pp. 132--147. \href{https://doi.org/10.1145/62212.62224}{DOI: 10.1145/62212.62224}.
\par\medskip

\noindent Wesley H. Holliday (2009). Dynamic testimonial logic. In Xiangdong He, John Horty and Eric Pacuit (eds.), Logic, Rationality, and Interaction, LORI 2009, Lecture Notes in Computer Science 5834, Springer, pp. 161--179. \href{https://doi.org/10.1007/978-3-642-04893-7_13}{DOI: 10.1007/978-3-642-04893-7\_13}. Trust put into a logic of changing your mind, with a graph of whose word each person takes.
\par\medskip

\noindent Radha Jagadeesan, Alan Jeffrey, Corin Pitcher and James Riely (2009). Towards a theory of accountability and audit. Computer Security, ESORICS 2009, Lecture Notes in Computer Science 5789, Springer, pp. 152--167. \href{https://doi.org/10.1007/978-3-642-04444-1_10}{DOI: 10.1007/978-3-642-04444-1\_10}. Systems that enforce their rules after the fact, by audit.
\par\medskip

\noindent Christoph J\"ager and Nicholas Shackel (2025). Testimonial authority and knowledge transmission. Social Epistemology 39(6), pp. 636--651. \href{https://doi.org/10.1080/02691728.2025.2449608}{DOI: 10.1080/02691728.2025.2449608}. A model where the speaker knowing is neither needed nor enough for the hearer to know.
\par\medskip

\noindent Rudolf E. Kalman (1960). A New Approach to Linear Filtering and Prediction Problems. Journal of Basic Engineering 82(1), pp. 35--45. \href{https://doi.org/10.1115/1.3662552}{DOI: 10.1115/1.3662552}. The filter.
\par\medskip

\noindent Rudolf E. Kalman (1960). On the general theory of control systems. Proceedings of the First IFAC World Congress, Moscow; IFAC Proceedings Volumes 1(1), pp. 491--502. \href{https://doi.org/10.1016/S1474-6670(17)70094-8}{DOI: 10.1016/S1474-6670(17)70094-8}. Observability, with a test you can compute.
\par\medskip

\noindent Oliver Kosut and J\"org Kliewer (2018). Authentication capacity of adversarial channels. 2018 IEEE Information Theory Workshop (ITW), Guangzhou, pp. 1--5. \href{https://doi.org/10.1109/ITW.2018.8613353}{DOI: 10.1109/ITW.2018.8613353}. The capacity is either zero or the full capacity without an attacker, depending on whether the attacker can make what's received look like an honest message.
\par\medskip

\noindent Robert K\"unnemann, Ilkan Esiyok and Michael Backes (2019). Automated verification of accountability in security protocols. 2019 IEEE 32nd Computer Security Foundations Symposium (CSF). \href{https://doi.org/10.1109/CSF.2019.00034}{DOI: 10.1109/CSF.2019.00034}; \href{https://arxiv.org/abs/1805.10891}{arXiv: 1805.10891}. Accountability checked automatically by making it a property of the whole run.
\par\medskip

\noindent Ralf K\"usters, Tomasz Truderung and Andreas Vogt (2010). Accountability: definition and relationship to verifiability. Proceedings of the 17th ACM Conference on Computer and Communications Security (CCS '10), pp. 526--535. \href{https://doi.org/10.1145/1866307.1866366}{DOI: 10.1145/1866307.1866366}. A verdict that names who misbehaved, with being able to check a result as a weaker form of it.
\par\medskip

\noindent Jennifer Lackey (2008). Learning from Words: Testimony as a Source of Knowledge. Oxford University Press. The teacher who doesn't believe the theory she teaches, and students who come away knowing it.
\par\medskip

\noindent F. William Lawvere (1969). Adjointness in foundations. Dialectica 23(3--4), pp. 281--296. \href{https://doi.org/10.1111/j.1746-8361.1969.tb01194.x}{DOI: 10.1111/j.1746-8361.1969.tb01194.x}. The two quantifiers as the two adjoints of substitution.
\par\medskip

\noindent William F. Lindgren and Thomas F. Banchoff (2009). Flatland: An Edition with Notes and Commentary, by Edwin A. Abbott. Cambridge University Press and Mathematical Association of America. \href{https://doi.org/10.1017/CBO9781139194921}{DOI: 10.1017/CBO9781139194921}. Calls the pen name a pun on Abbott's initials, EAA.
\par\medskip

\noindent John Locke (1690). An Essay Concerning Human Understanding, Book I, chapter iv, section 23. \href{https://www.earlymoderntexts.com/assets/pdfs/locke1690book1.pdf}{Modernized text of Book I}. Believing something on authority isn't knowing it.
\par\medskip

\noindent Leonard A. McGee and Stanley F. Schmidt (1985). Discovery of the Kalman Filter as a Practical Tool for Aerospace and Industry. NASA Technical Memorandum 86847, Ames Research Center. \href{https://ntrs.nasa.gov/citations/19860003843}{NTRS}. Kalman's visit to Ames in the fall of 1960 and the filter's road into navigation for Apollo.
\par\medskip

\noindent Edward F. Moore (1956). Gedanken-Experiments on Sequential Machines. In Claude E. Shannon and John McCarthy (eds.), Automata Studies, Princeton University Press, pp. 129--153. \href{https://www.cs.cmu.edu/~cdm/resources/Moore1956-gedanken-experiments.pdf}{Original paper}.
\par\medskip

\noindent David L. Parnas (1972). On the criteria to be used in decomposing systems into modules. Communications of the ACM 15(12), pp. 1053--1058. \href{https://doi.org/10.1145/361598.361623}{DOI: 10.1145/361598.361623}. A module as ``a responsibility assignment'' on p.~1054, and the decomposition that could be ``much less efficient'' on p.~1057.
\par\medskip

\noindent Bryan Parno, Jon Howell, Craig Gentry and Mariana Raykova (2013). Pinocchio: nearly practical verifiable computation. Proceedings of the 2013 IEEE Symposium on Security and Privacy, pp. 238--252. \href{https://doi.org/10.1109/SP.2013.47}{DOI: 10.1109/SP.2013.47}.
\par\medskip

\noindent Jean-Jacques Quisquater, Myriam Quisquater, Muriel Quisquater, Micha\"el Quisquater, Louis Guillou, Marie Annick Guillou, Ga\"id Guillou, Anna Guillou, Gwenol\'e Guillou, Soazig Guillou and Thomas A. Berson (1990). How to Explain Zero-Knowledge Protocols to Your Children. Advances in Cryptology, CRYPTO '89, Lecture Notes in Computer Science 435, Springer, pp. 628--631. \href{https://doi.org/10.1007/0-387-34805-0_60}{DOI: 10.1007/0-387-34805-0\_60}. Ali Baba's cave, told for children.
\par\medskip

\noindent Plato. Republic, Book VII, 514a--517c, translated by Benjamin Jowett. \href{https://classics.mit.edu/Plato/republic.8.vii.html}{MIT Internet Classics Archive}. The cave; the prisoners' honours for whoever is quickest at the passing shadows and best at guessing what comes next (516c--d); and Socrates' own reading of the cave as the world of sight (517b).
\par\medskip

\noindent Henry Gordon Rice (1953). Classes of Recursively Enumerable Sets and Their Decision Problems. Transactions of the American Mathematical Society 74(2), pp. 358--366. \href{https://doi.org/10.1090/S0002-9947-1953-0053041-6}{DOI: 10.1090/S0002-9947-1953-0053041-6}.
\par\medskip

\noindent Rudy Rucker (1984). The Fourth Dimension: A Guided Tour of the Higher Universes. Houghton Mifflin, Boston. The earliest guess found that ``A Square'' was a nickname for A squared.
\par\medskip

\noindent Anooshiravan Saboori and Christoforos N. Hadjicostis (2007). Notions of security and opacity in discrete event systems. Proceedings of the 46th IEEE Conference on Decision and Control, pp. 5056--5061. \href{https://doi.org/10.1109/CDC.2007.4434515}{DOI: 10.1109/CDC.2007.4434515}. When an observer can never be sure the system is in a secret state.
\par\medskip

\noindent Meera Sampath, Raja Sengupta, St\'ephane Lafortune, Kasim Sinnamohideen and Demosthenis Teneketzis (1995). Diagnosability of discrete-event systems. IEEE Transactions on Automatic Control 40(9), pp. 1555--1575. \href{https://doi.org/10.1109/9.412626}{DOI: 10.1109/9.412626}. When an observer can always work out from what it sees that a hidden fault has happened.
\par\medskip

\noindent Bruce Schneier and John Kelsey (1999). Secure audit logs to support computer forensics. ACM Transactions on Information and System Security 2(2), pp. 159--176. \href{https://doi.org/10.1145/317087.317089}{DOI: 10.1145/317087.317089}. A log an intruder who takes over the machine can't quietly rewrite, protecting what was written before the break-in and nothing after.
\par\medskip

\noindent Ian Stewart (2002). The Annotated Flatland: A Romance of Many Dimensions. Perseus Publishing. Abbott's text with Stewart's introduction and notes.
\par\medskip

\noindent Anvith Thudi, Hengrui Jia, Ilia Shumailov and Nicolas Papernot (2022). On the Necessity of Auditable Algorithmic Definitions for Machine Unlearning. Proceedings of the 31st USENIX Security Symposium. \href{https://arxiv.org/abs/2110.11891}{arXiv: 2110.11891}. The same finished model can come out of different training sets, so only the procedure that was run can be checked.
\par\medskip

\noindent Timo Treitz and Robert K\"unnemann (2026). Accountability in Certificate Transparency and variants. Preprint. \href{https://arxiv.org/abs/2609.11552}{arXiv: 2609.11552}. Plain Certificate Transparency names a culprit only if the log itself is honest.
\par\medskip

\noindent Andreas Winter, Anderson C. A. Nascimento and Hideki Imai (2003). Commitment capacity of discrete memoryless channels. Cryptography and Coding 2003, Lecture Notes in Computer Science 2898, Springer, pp. 35--51. \href{https://doi.org/10.1007/978-3-540-40974-8_4}{DOI: 10.1007/978-3-540-40974-8\_4}; \href{https://arxiv.org/abs/cs/0304014}{arXiv: cs/0304014}. How fast a noisy channel can carry commitments.
\par\medskip

\subsection*{Chapter 8. Established}

\noindent Federal Aviation Administration (2026). 25-Hour Cockpit Voice Recorder (CVR) Requirement, New Aircraft Production. Final rule, Federal Register 91, p. 4447, 2 February 2026; docket FAA-2023-2270. \href{https://www.govinfo.gov/content/pkg/FR-2026-02-02/html/2026-02110.htm}{Federal Register text}. Raises the required recording time from two hours to twenty-five for newly manufactured aircraft.
\par\medskip

\noindent Arthur Azevedo de Amorim, Maxime D\'en\`es, Nick Giannarakis, C\u{a}t\u{a}lin Hri\c{t}cu, Benjamin C. Pierce, Antal Spector-Zabusky and Andrew Tolmach (2015). Micro-policies: formally verified, tag-based security monitors. Proceedings of the 2015 IEEE Symposium on Security and Privacy, pp. 813--830. \href{https://doi.org/10.1109/SP.2015.55}{DOI: 10.1109/SP.2015.55}. Monitors checked at every instruction, with their guarantees proved.
\par\medskip

\noindent James P. Anderson (1972). Computer Security Technology Planning Study, Volume II. Technical Report ESD-TR-73-51, Electronic Systems Division, Air Force Systems Command, Hanscom Field, Bedford, MA. \href{https://csrc.nist.gov/csrc/media/publications/conference-paper/1998/10/08/proceedings-of-the-21st-nissc-1998/documents/early-cs-papers/ande72.pdf}{Report}. The reference validation mechanism ``must be tamper proof,'' ``must always be invoked,'' and ``must be small enough to be subject to analysis and tests.''
\par\medskip

\noindent Luis Antunes, Lance Fortnow, Dieter van Melkebeek and N. V. Vinodchandran (2006). Computational depth: Concept and applications. Theoretical Computer Science 354(3), pp. 391--404. \href{https://doi.org/10.1016/j.tcs.2005.11.033}{DOI: 10.1016/j.tcs.2005.11.033}. Its epigraph is Levin's line: ``Torturing an uninformed witness cannot give information about the crime.''
\par\medskip

\noindent Sergei N. Artemov (1995). Operational Modal Logic. Technical Report MSI 95-29, Cornell University. The logic of proofs, where knowing a thing comes with an explicit reason for it; surveyed in the Stanford Encyclopedia of Philosophy entry on justification logic, \href{https://plato.stanford.edu/entries/logic-justification/}{Entry}.
\par\medskip

\noindent Sergei Artemov and Roman Kuznets (2014). Logical omniscience as infeasibility. Annals of Pure and Applied Logic 165(1), pp. 6--25. \href{https://doi.org/10.1016/j.apal.2013.07.003}{DOI: 10.1016/j.apal.2013.07.003}. Knowing that counts the steps it took.
\par\medskip

\noindent J. L. Austin (1962). How to Do Things with Words. Clarendon Press, Oxford. The William James Lectures, Harvard, 1955, where saying ``I do'' makes a marriage rather than describing one.
\par\medskip

\noindent Francisco Bavera and Eduardo Bonelli (2018). Justification logic and audited computation. Journal of Logic and Computation 28(5), pp. 909--934; first published online in 2015. \href{https://doi.org/10.1093/logcom/exv037}{DOI: 10.1093/logcom/exv037}. A programming language whose terms carry the history of the computation that made them.
\par\medskip

\noindent Matt Bishop and Michael Dilger (1996). Checking for race conditions in file accesses. Computing Systems 9(2), pp. 131--152. \href{https://www.usenix.org/legacy/publications/compsystems/1996/spr_bishop.pdf}{PDF}. Hunting for time-of-check to time-of-use bugs.
\par\medskip

\noindent Samuel Bucheli, Meghdad Ghari and Thomas Studer (2017). Temporal justification logic. Methods for Modalities (M4M9), Electronic Proceedings in Theoretical Computer Science 243, pp. 59--74. \href{https://doi.org/10.4204/EPTCS.243.5}{DOI: 10.4204/EPTCS.243.5}; \href{https://arxiv.org/abs/1703.02195}{arXiv: 1703.02195}.
\par\medskip

\noindent Eric Budish (2025). Trust at Scale: The Economic Limits of Cryptocurrencies and Blockchains. The Quarterly Journal of Economics 140(1), pp. 1--62. \href{https://doi.org/10.1093/qje/qjae033}{DOI: 10.1093/qje/qjae033}. For Bitcoin's kind of trust to hold, honest work has to stay expensive at every moment compared with what an attacker could win.
\par\medskip

\noindent Gregory J. Chaitin (1966). On the Length of Programs for Computing Finite Binary Sequences. Journal of the ACM 13(4), pp. 547--569. \href{https://doi.org/10.1145/321356.321363}{DOI: 10.1145/321356.321363}.
\par\medskip

\noindent Ho Ngoc Duc (1997). Reasoning about rational, but not logically omniscient, agents. Journal of Logic and Computation 7(5), pp. 633--648. \href{https://doi.org/10.1093/logcom/7.5.633}{DOI: 10.1093/logcom/7.5.633}. Knowing as what an agent could come to know with more steps of reasoning.
\par\medskip

\noindent Edmund L. Gettier (1963). Is Justified True Belief Knowledge? Analysis 23(6), pp. 121--123. \href{https://doi.org/10.1093/analys/23.6.121}{DOI: 10.1093/analys/23.6.121}. Smith, Jones and the ten coins.
\par\medskip

\noindent Alvin I. Goldman (1976). Discrimination and Perceptual Knowledge. The Journal of Philosophy 73(20), pp. 771--791. \href{https://doi.org/10.2307/2025679}{DOI: 10.2307/2025679}. The fake-barn case, credited there to Carl Ginet.
\par\medskip

\noindent Michael J. C. Gordon, Robin Milner and Christopher P. Wadsworth (1979). Edinburgh LCF: A Mechanised Logic of Computation. Lecture Notes in Computer Science 78, Springer. \href{https://doi.org/10.1007/3-540-09724-4}{DOI: 10.1007/3-540-09724-4}.
\par\medskip

\noindent Michael J. C. Gordon (2000). From LCF to HOL: a short history. In Gordon Plotkin, Colin Stirling and Mads Tofte (eds.), Proof, Language, and Interaction: Essays in Honour of Robin Milner, MIT Press, pp. 169--185. \href{https://www.cl.cam.ac.uk/archive/mjcg/papers/HolHistory.pdf}{Paper}. Milner keeping only the theorems, ``like a mathematics lecturer using a small blackboard,'' with an abstract type so that theorems can only come from proofs, and ML named for ``Meta Language.''
\par\medskip

\noindent Jim Gray (1981). The Transaction Concept: Virtues and Limitations. Proceedings of the 7th International Conference on Very Large Data Bases, pp. 144--154. \href{https://jimgray.azurewebsites.net/papers/thetransactionconcept.pdf}{Paper}. The commit, introduced through contract law and a wedding.
\par\medskip

\noindent Trusted Computing Group (n.d.). Trusted Platform Module Library Specification, Family 2.0, Part 1: Architecture. \href{https://trustedcomputinggroup.org/resource/tpm-library-specification/}{Specification page}. The registers and the extend operation, the bit-field slots that can only set bits, and the counters that only go up.
\par\medskip

\noindent Jaakko Hintikka (1962). Knowledge and Belief: An Introduction to the Logic of the Two Notions. Cornell University Press. Logical omniscience, surveyed in the Stanford Encyclopedia of Philosophy entry on epistemic logic, \href{https://plato.stanford.edu/entries/logic-epistemic/}{Entry}.
\par\medskip

\noindent Andrei N. Kolmogorov (1965). Three approaches to the quantitative definition of information. Problems of Information Transmission 1(1), pp. 1--7; reprinted in International Journal of Computer Mathematics 2 (1968), pp. 157--168. \href{https://doi.org/10.1080/00207166808803030}{DOI of the reprint: 10.1080/00207166808803030}.
\par\medskip

\noindent H. T. Kung and John T. Robinson (1981). On optimistic methods for concurrency control. ACM Transactions on Database Systems 6(2), pp. 213--226. \href{https://doi.org/10.1145/319566.319567}{DOI: 10.1145/319566.319567}. Checking, just before a commit, that nothing read has been written since.
\par\medskip

\noindent Leonid A. Levin (1984). Randomness Conservation Inequalities; Information and Independence in Mathematical Theories. Information and Control 61(1), pp. 15--37. \href{https://doi.org/10.1016/S0019-9958(84)80060-1}{DOI: 10.1016/S0019-9958(84)80060-1}. Running a program can't add much information beyond what went in plus the program.
\par\medskip

\noindent Jacques-Louis Lions (chair) et al. (1996). ARIANE 5 Flight 501 Failure: Report by the Inquiry Board. Paris, 19 July 1996. \href{https://www-users.cse.umn.edu/~arnold/disasters/ariane5rep.html}{Copy of the report}. The assumption that made the conversion safe was agreed on and then ``essentially obscured, though not deliberately, from any external review.''
\par\medskip

\noindent Moni Naor and Guy N. Rothblum (2005). The complexity of online memory checking. Proceedings of the 46th IEEE Symposium on Foundations of Computer Science (FOCS 2005); journal version in Journal of the ACM 56(1), 2009. \href{https://doi.org/10.1145/1462153.1462155}{DOI of the journal version: 10.1145/1462153.1462155}; \href{https://eprint.iacr.org/2006/091}{IACR ePrint 2006/091}. Against an unbounded opponent, the size of the client's private fingerprint times the bits it reads has to be on the order of the whole file, and beating that needs one-way functions.
\par\medskip

\noindent National Museum of Australia. Defining Moments: Black box flight recorder invented. \href{https://www.nma.gov.au/defining-moments/resources/black-box-invented}{Museum page}. David Warren's flight memory recorder at the Aeronautical Research Laboratories in Melbourne.
\par\medskip

\noindent Kerry Noble (2013). Obituary: Dr David Ronald De Mey Warren (1925--2010). Imperial College London News, 14 June 2013. \href{https://www.imperial.ac.uk/news/122905/obituary-dr-david-ronald-de-mey/}{Article}. Warren at the Aeronautical Research Laboratories, the unsolved Comet crashes, and the 1956 prototype that held four hours of voice and instrument data.
\par\medskip

\noindent Robert Nozick (1981). Philosophical Explanations. Harvard University Press. Chapter 3: you know a thing only if, had it been false, you wouldn't believe it.
\par\medskip

\noindent Daniel Oberhaus (2018). The World's Oldest Blockchain Has Been Hiding in the New York Times Since 1995. Vice (Motherboard), 27 August 2018. \href{https://www.vice.com/en/article/what-was-the-first-blockchain/}{Article}. Surety's weekly hash under ``Notices \& Lost and Found.''
\par\medskip

\noindent Plato. Meno 97d--98a, translated by W. R. M. Lamb, Plato in Twelve Volumes, vol. 3, Harvard University Press, 1967. \href{http://www.perseus.tufts.edu/hopper/text?doc=Plat.+Meno+97d}{Perseus text}. Daedalus's statues, and true opinions tied down by reasoning.
\par\medskip

\noindent Fred B. Schneider (2000). Enforceable security policies. ACM Transactions on Information and System Security 3(1), pp. 30--50. \href{https://doi.org/10.1145/353323.353382}{DOI: 10.1145/353323.353382}. A monitor that watches a run as it goes can enforce exactly the rules whose breaking shows up at a point you can name.
\par\medskip

\noindent Claude E. Shannon (1948). A mathematical theory of communication. Bell System Technical Journal 27(3), pp. 379--423, and 27(4), pp. 623--656. \href{https://doi.org/10.1002/j.1538-7305.1948.tb01338.x}{DOI: 10.1002/j.1538-7305.1948.tb01338.x}; \href{https://doi.org/10.1002/j.1538-7305.1948.tb00917.x}{DOI: 10.1002/j.1538-7305.1948.tb00917.x}. Structure ``allows one to make a saving in time (or channel capacity),'' p.~385; half of English ``determined by the structure of the language,'' p.~399.
\par\medskip

\noindent Ray J. Solomonoff (1964). A Formal Theory of Inductive Inference, Parts I and II. Information and Control 7(1), pp. 1--22, and 7(2), pp. 224--254. \href{https://doi.org/10.1016/S0019-9958(64)90223-2}{DOI of Part I: 10.1016/S0019-9958(64)90223-2}; \href{https://doi.org/10.1016/S0019-9958(64)90131-7}{DOI of Part II: 10.1016/S0019-9958(64)90131-7}.
\par\medskip

\subsection*{Chapter 9. Settled}

\noindent Hans Albert (1968). Traktat \"uber kritische Vernunft. J. C. B. Mohr (Paul Siebeck), T\"ubingen. English translation: Treatise on Critical Reason, Princeton University Press, 1985. The M\"unchhausen trilemma.
\par\medskip

\noindent Kenneth Appel and Wolfgang Haken (1976). Every planar map is four colorable. Bulletin of the American Mathematical Society 82(5), pp. 711--712. \href{https://doi.org/10.1090/S0002-9904-1976-14122-5}{DOI: 10.1090/S0002-9904-1976-14122-5}. The announcement of the computer proof.
\par\medskip

\noindent Mark van Atten (2018). The Creating Subject, the Brouwer--Kripke Schema, and infinite proofs. Indagationes Mathematicae 29(6), pp. 1565--1636. \href{https://doi.org/10.1016/j.indag.2018.06.005}{DOI: 10.1016/j.indag.2018.06.005}; \href{https://arxiv.org/abs/1805.00404}{arXiv: 1805.00404}. Read in the arXiv preprint: the never-forgetting axiom credited to Troelstra (1969, p.~95), the name Kreisel-Troelstra theory, the reading ``by stage n the Creating Subject has made A evident,'' and the subject who can't simply accept a proof from outside.
\par\medskip

\noindent L. E. J. Brouwer (1948). Essentieel-negatieve eigenschappen. Proceedings of the Koninklijke Nederlandse Akademie van Wetenschappen 51, pp. 963--964; also Indagationes Mathematicae 10. English translation, Essentially negative properties, in Hans Freudenthal (ed.), L. E. J. Brouwer: Collected Works, vol.~2, North-Holland, 1976, pp. 478--479. Where the creating subject first appears in an argument.
\par\medskip

\noindent Walter Dean and Hidenori Kurokawa (2021). On the methodology of informal rigour: set theory, semantics, and intuitionism. Preprint, written for Marianna Antonutti Marfori and Mattia Petrolo (eds.), Intuitionism, Computation, and Proof: Selected Themes from the Research of G. Kreisel, Springer. \href{https://arxiv.org/abs/2104.14887}{arXiv: 2104.14887}. Its list of Kreisel's three creating-subject axioms is decidability, Christian charity and soundness, with no never-forgetting axiom.
\par\medskip

\noindent FIDE (2023). FIDE Laws of Chess taking effect from 1 January 2023, Article 9.2.3. \href{https://handbook.fide.com/chapter/E012023}{FIDE Handbook}. Two positions are the same only if the same player has the move, the same pieces stand on the same squares, and the possible moves of all the pieces are the same.
\par\medskip

\noindent Carl Friedrich Gauss (1831). Untersuchungen \"uber die Eigenschaften der positiven tern\"aren quadratischen Formen von Ludwig August Seeber. G\"ottingische gelehrte Anzeigen, 9 July 1831; reprinted in Journal f\"ur die reine und angewandte Mathematik 20 (1840), pp. 312--320. \href{https://doi.org/10.1515/crll.1840.20.312}{DOI of the reprint: 10.1515/crll.1840.20.312}. The densest packing among lattices, a special case of Kepler's claim.
\par\medskip

\noindent Georges Gonthier (2008). Formal proof: the four-color theorem. Notices of the American Mathematical Society 55(11), pp. 1382--1393. The proof checked step by step by a small proof checker, completed in 2005.
\par\medskip

\noindent Thomas C. Hales (2005). A proof of the Kepler conjecture. Annals of Mathematics 162(3), pp. 1065--1185. \href{https://doi.org/10.4007/annals.2005.162.1065}{DOI: 10.4007/annals.2005.162.1065}.
\par\medskip

\noindent Thomas Hales and others (2017). A formal proof of the Kepler conjecture. Forum of Mathematics, Pi 5, e2. \href{https://doi.org/10.1017/fmp.2017.1}{DOI: 10.1017/fmp.2017.1}. The project announced in 2003 and completed on 10 August 2014.
\par\medskip

\noindent Thomas Hales (2024). The Formal Proof of the Kepler Conjecture: a critical retrospective. \href{https://arxiv.org/abs/2402.08032}{arXiv: 2402.08032}. The August 1998 announcement with Ferguson, the panel of twelve referees who ``became exhausted and quit,'' and the formal proof begun to get past them (p.~2).
\par\medskip

\noindent Percy John Heawood (1890). Map-colour theorem. Quarterly Journal of Pure and Applied Mathematics 24, pp. 332--338. The flaw in ``the now apparently recognised proof,'' as quoted in secondary accounts; I haven't seen the original.
\par\medskip

\noindent MacTutor History of Mathematics. Percy John Heawood (J. J. O'Connor and E. F. Robertson). University of St Andrews. \href{https://mathshistory.st-andrews.ac.uk/Biographies/Heawood/}{Biography}. Lecturer at Durham from 1887, and, in G. A. Dirac's words, ``often accompanied by a dog, which was admitted to his lectures.''
\par\medskip

\noindent David Hilbert (1900). Mathematische Probleme. Nachrichten von der K\"oniglichen Gesellschaft der Wissenschaften zu G\"ottingen, pp. 253--297. English translation by Mary Winston Newson, Mathematical Problems, Bulletin of the American Mathematical Society 8 (1902), pp. 437--479, \href{https://doi.org/10.1090/S0002-9904-1902-00923-3}{DOI: 10.1090/S0002-9904-1902-00923-3}. The packing question sits in the eighteenth problem.
\par\medskip

\noindent Canterbury Historical and Archaeological Society (n.d.). James Six (1730--1793). \href{https://www.canterbury-archaeology.org.uk/six-james-biog}{Web page}. The family silk-weaving business and his turn to astronomy and meteorology.
\par\medskip

\noindent Rosalie Iemhoff (2008, revised 2019). Intuitionism in the Philosophy of Mathematics. Stanford Encyclopedia of Philosophy. \href{https://plato.stanford.edu/entries/intuitionism/}{Entry}. Lists ``the Creating Subject never forgets what it has proven'' among the three axioms Kreisel introduced in 1967.
\par\medskip

\noindent Alfred Bray Kempe (1879). On the Geographical Problem of the Four Colours. American Journal of Mathematics 2(3), pp. 193--200. \href{https://doi.org/10.2307/2369235}{DOI: 10.2307/2369235}. The proof that stood for eleven years; his account of finding it is quoted from p.~194.
\par\medskip

\noindent MacTutor History of Mathematics. Alfred Bray Kempe (J. J. O'Connor and E. F. Robertson). University of St Andrews. \href{https://mathshistory.st-andrews.ac.uk/Biographies/Kempe/}{Biography}. A barrister with ``a very fine counter-tenor voice.''
\par\medskip

\noindent Johannes Kepler (1611). Strena seu de Nive Sexangula. Godefrid Tampach, Frankfurt am Main. English: The Six-Cornered Snowflake, translated by Colin Hardie, with essays by L. L. Whyte and B. J. Mason, Clarendon Press, Oxford, 1966. The New Year's gift to Wackher, and the claim about stacking balls; I haven't seen the Hardie translation.
\par\medskip

\noindent Lewis A. Kornhauser and Lawrence G. Sager (1986). Unpacking the court. Yale Law Journal 96(1), beginning p. 82. \href{https://doi.org/10.2307/796436}{DOI: 10.2307/796436}. A court that votes by majority on each issue and on the verdict can contradict itself.
\par\medskip

\noindent Georg Kreisel (1967). Informal rigour and completeness proofs. In Imre Lakatos (ed.), Problems in the Philosophy of Mathematics, Studies in Logic and the Foundations of Mathematics 47, North-Holland, Amsterdam, pp. 138--186. \href{https://doi.org/10.1016/S0049-237X(08)71525-8}{DOI: 10.1016/S0049-237X(08)71525-8}. The creating subject's evidence for a statement at stage n.
\par\medskip

\noindent Jeffrey C. Lagarias (ed.) (2011). The Kepler Conjecture: The Hales--Ferguson Proof. Springer, New York. \href{https://doi.org/10.1007/978-1-4614-1129-1}{DOI: 10.1007/978-1-4614-1129-1}. Reprints the six papers by Thomas C. Hales and Samuel P. Ferguson that gave the full proof in Discrete \& Computational Geometry 36(1), 2006.
\par\medskip

\noindent Christian List and Philip Pettit (2002). Aggregating sets of judgments: an impossibility result. Economics and Philosophy 18(1), pp. 89--110. \href{https://doi.org/10.1017/S0266267102001098}{DOI: 10.1017/S0266267102001098}.
\par\medskip

\noindent Institute for Logic, Language and Computation (2019). Anne Troelstra (1939--2019). Obituary by Johan van Benthem and Dick de Jongh, University of Amsterdam. \href{https://www.illc.uva.nl/NewsandEvents/News/Obituaries/Anne-Troelstra-1939-2019-/}{Page}. His doctorate was supervised by Heyting.
\par\medskip

\noindent MacTutor History of Mathematics. Arend Heyting. University of St Andrews. \href{https://mathshistory.st-andrews.ac.uk/Biographies/Heyting/}{Biography}. Doctorate in 1925 under Brouwer.
\par\medskip

\noindent Per Martin-L\"of (1996). On the Meanings of the Logical Constants and the Justifications of the Logical Laws. Nordic Journal of Philosophical Logic 1(1), pp. 11--60. Lectures given in Siena, 6--9 April 1983. \href{https://ncatlab.org/nlab/files/MartinLofOnTheMeaning96.pdf}{Paper}. ``Thus a proof is, not an object, but an act,'' and ``Only secondarily, and irrevocably, does it become the act that has been performed.''
\par\medskip

\noindent Bureau International des Poids et Mesures (2019). The International System of Units (SI), 9th edition. \href{https://www.bipm.org/en/publications/si-brochure}{SI Brochure}. The kilogram defined by fixing the Planck constant, adopted by the 26th CGPM in 2018 and in force from 20 May 2019.
\par\medskip

\noindent Bureau International des Poids et Mesures (n.d.). Historical perspective: unit of mass, kilogram. \href{https://www.bipm.org/en/history-si/kilogram}{BIPM page}. The national copies found drifting from the international prototype in 1946, confirmed by the verification of 1989 to 1991, with a median difference of about 25 micrograms.
\par\medskip

\noindent Joint Committee for Guides in Metrology (2012). International vocabulary of metrology: Basic and general concepts and associated terms (VIM), 3rd edition. JCGM 200:2012. \href{https://www.bipm.org/documents/20126/2071204/JCGM_200_2012.pdf}{PDF}. Definition 2.41, metrological traceability ``through a documented unbroken chain of calibrations,'' and its Note~5, that traceability ``does not ensure \ldots{} that there is an absence of mistakes.''
\par\medskip

\noindent World Health Organization and UNICEF (2007). WHO and UNICEF urge donors and member states to fund and expand use of vaccine vial monitors to ensure effective immunization. Joint press release, Geneva, 3 May 2007. \href{https://www.gavi.org/news/media-room/who-and-unicef-urge-donors-and-member-states-fund-and-expand-use-vaccine-vial}{Release}. The first monitors on oral polio vaccine in 1996, with a heat-sensitive inner square that darkens as the vial takes on heat.
\par\medskip

\noindent Charles S. Peirce (1877). The Fixation of Belief. Popular Science Monthly 12, pp. 1--15. ``As soon as a firm belief is reached we are entirely satisfied, whether the belief be true or false,'' p.~6.
\par\medskip

\noindent Plato. Republic, Book VII, 514a--517c, translated by Benjamin Jowett. \href{https://classics.mit.edu/Plato/republic.8.vii.html}{MIT Internet Classics Archive}. The prisoners in the cave.
\par\medskip

\noindent Karl R. Popper (1959). The Logic of Scientific Discovery. Hutchinson, London. Translation of Logik der Forschung (1934). Chapter X, corroboration: credit earned by surviving a test that could have failed.
\par\medskip

\noindent Neil Robertson, Daniel Sanders, Paul Seymour and Robin Thomas (1997). The Four-Colour Theorem. Journal of Combinatorial Theory, Series B 70(1), pp. 2--44. \href{https://doi.org/10.1006/jctb.1997.1750}{DOI: 10.1006/jctb.1997.1750}. The cleaner computer proof.
\par\medskip

\noindent Sextus Empiricus. Outlines of Pyrrhonism I.164--177. The five modes credited to Agrippa.
\par\medskip

\noindent James Six (1782). Account of an improved Thermometer. Philosophical Transactions of the Royal Society of London 72, pp. 72--81. \href{https://doi.org/10.1098/rstl.1782.0012}{DOI: 10.1098/rstl.1782.0012}. ``If it be not immediately noticed, it can never after be known,'' p.~72.
\par\medskip

\noindent William E. Story (1879). Note on the Preceding Paper. American Journal of Mathematics 2(3), from p. 201. \href{https://doi.org/10.2307/2369236}{DOI: 10.2307/2369236}. ``Substantially proved,'' p.~201.
\par\medskip

\noindent G\"oran Sundholm (2014). Constructive Recursive Functions, Church's Thesis, and Brouwer's Theory of the Creating Subject: Afterthoughts on a Parisian Joint Session. In Jacques Dubucs and Michel Bourdeau (eds.), Constructivity and Computability in Historical and Philosophical Perspective, Logic, Epistemology, and the Unity of Science 34, Springer, pp. 1--35. \href{https://doi.org/10.1007/978-94-017-9217-2_1}{DOI: 10.1007/978-94-017-9217-2\_1}.
\par\medskip

\noindent George G. Szpiro (2003). Does the proof stack up? Nature 424(6944), pp. 12--13. \href{https://doi.org/10.1038/424012a}{DOI: 10.1038/424012a}. The Annals editor's letter as Hales reported it, ending ``run out of energy.''
\par\medskip

\noindent MacTutor History of Mathematics. The four colour theorem (J. J. O'Connor and E. F. Robertson). University of St Andrews. \href{https://mathshistory.st-andrews.ac.uk/HistTopics/The_four_colour_theorem/}{History topic}. De Morgan's letter to Hamilton of 23 October 1852, asking for ``a reason for a fact which I did not know was a fact.''
\par\medskip

\noindent L\'aszl\'o Fejes T\'oth (1953). Lagerungen in der Ebene, auf der Kugel und im Raum. Springer, Berlin. The first detailed strategy for proving Kepler's claim, and the first suggestion that a computer might finish it.
\par\medskip

\noindent Anne S. Troelstra (1969). Principles of Intuitionism. Lecture Notes in Mathematics 95, Springer, Berlin. \href{https://doi.org/10.1007/BFb0080643}{DOI: 10.1007/BFb0080643}. The creating subject who never forgets, added on p.~95.
\par\medskip

\noindent Robin Wilson (2002). Four Colours Suffice: How the Map Problem Was Solved. Allen Lane, London. Francis Guthrie and the map of the counties of England.
\par\medskip

\subsection*{Chapter 10. What structure is}

\noindent Steve Awodey (1996). Structure in mathematics and logic: a categorical perspective. Philosophia Mathematica 4(3), pp. 209--237. \href{https://doi.org/10.1093/philmat/4.3.209}{DOI: 10.1093/philmat/4.3.209}.
\par\medskip

\noindent Michael J. Barany (2020). Impersonation and personification in mid-twentieth century mathematics. History of Science 58(4), pp. 417--436. \href{https://doi.org/10.1177/0073275320924571}{DOI: 10.1177/0073275320924571}. Bourbaki's applications to the American Mathematical Society in 1948 and again in 1949, and the letters around them.
\par\medskip

\noindent Paul Benacerraf (1965). What numbers could not be. The Philosophical Review 74(1), pp. 47--73. \href{https://doi.org/10.2307/2183530}{DOI: 10.2307/2183530}. Two equally good ways of building the numbers out of sets that disagree about which set three is.
\par\medskip

\noindent Antoine B\'erut, Artak Arakelyan, Artyom Petrosyan, Sergio Ciliberto, Raoul Dillenschneider and Eric Lutz (2012). Experimental verification of Landauer's principle linking information and thermodynamics. Nature 483, pp. 187--189. \href{https://doi.org/10.1038/nature10872}{DOI: 10.1038/nature10872}. One bit erased in a single colloidal bead at the \'Ecole Normale Sup\'erieure in Lyon, the mean heat approaching Landauer's bound as the erasure is slowed.
\par\medskip

\noindent J\'anos Bolyai (1832). Appendix scientiam spatii absolute veram exhibens. Printed as an appendix to the first volume of Farkas Bolyai's Tentamen, Maros-V\'as\'arhely. The son's geometry without the parallel postulate.
\par\medskip

\noindent Roberto Bonola (1912). Non-Euclidean Geometry: A Critical and Historical Study of Its Development. Translated by H. S. Carslaw. Open Court, Chicago. \href{https://archive.org/details/noneuclideangeom00bonouoft}{Internet Archive}. Two thousand years of attempts on the parallel postulate, from Proclus and the Arab geometers through Saccheri to Gauss, Bolyai and Lobachevsky.
\par\medskip

\noindent George Boole (1847). The Mathematical Analysis of Logic, Being an Essay towards a Calculus of Deductive Reasoning. Macmillan, Barclay and Macmillan, Cambridge.
\par\medskip

\noindent George Boole (1854). An Investigation of the Laws of Thought. Walton and Maberly, London.
\par\medskip

\noindent Nicolas Bourbaki (1950). The Architecture of Mathematics. Translated by Arnold Dresden. American Mathematical Monthly 57(4), pp. 221--232. \href{https://doi.org/10.1080/00029890.1950.11999523}{DOI: 10.1080/00029890.1950.11999523}. Structures on sets of elements ``whose nature has not been specified,'' and the three mother structures, algebraic, order and topological.
\par\medskip

\noindent Nicolas Bourbaki (1968; reprinted 2004). Theory of Sets. Elements of Mathematics, Hermann, Paris, and Addison-Wesley; Springer reprint, Berlin, 2004. \href{https://doi.org/10.1007/978-3-642-59309-3}{DOI: 10.1007/978-3-642-59309-3}. Chapter IV, sections 2 and 3, on morphisms, derived structures and universal mappings.
\par\medskip

\noindent Walter Burkert (1972). Lore and Science in Ancient Pythagoreanism. Translated by Edwin L. Minar Jr. Harvard University Press, Cambridge, Massachusetts. Read here through Huffman's entry on Pythagoreanism.
\par\medskip

\noindent Leo Corry (1992). Nicolas Bourbaki and the concept of mathematical structure. Synthese 92(3), pp. 315--348. \href{https://doi.org/10.1007/BF00414286}{DOI: 10.1007/BF00414286}. The promised chapter on categories that never appeared in the treatise.
\par\medskip

\noindent Leo Corry (1996; 2nd edition 2004). Modern Algebra and the Rise of Mathematical Structures. Birkh\"auser, Basel. \href{https://doi.org/10.1007/978-3-0348-7917-0}{DOI: 10.1007/978-3-0348-7917-0}. Chapter 7, on Bourbaki's theory of structures; ``plays no actual role in the presentation of theories within the treatise'' is on p.~338.
\par\medskip

\noindent Samuel Eilenberg and Saunders Mac Lane (1945). General Theory of Natural Equivalences. Transactions of the American Mathematical Society 58(2), pp. 231--294. \href{https://doi.org/10.1090/S0002-9947-1945-0013131-6}{DOI: 10.1090/S0002-9947-1945-0013131-6}. Where category theory starts: categories, functors and natural transformations.
\par\medskip

\noindent Euclid. Elements. Translated by T. L. Heath, The Thirteen Books of Euclid's Elements, Cambridge University Press, 1908. \href{http://www.perseus.tufts.edu/hopper/text?doc=Perseus:text:1999.01.0086}{Perseus text}. Definition 1, Postulate 1, the fifth postulate, Common Notions 1 and 4, and Proposition 4 of Book I, proved by laying one triangle on another.
\par\medskip

\noindent Gottlob Frege (1884). Die Grundlagen der Arithmetik. Wilhelm Koebner, Breslau. \href{https://www.gutenberg.org/ebooks/48312}{Project Gutenberg 48312}. Directions and parallel lines in section 64, and number through one-to-one correspondence in sections 68 to 73.
\par\medskip

\noindent Gottlob Frege (1893). Grundgesetze der Arithmetik, vol. I. Hermann Pohle, Jena. Basic Law V, on value-ranges, in section 20.
\par\medskip

\noindent Gottlob Frege (1903). Grundgesetze der Arithmetik, vol. II. Hermann Pohle, Jena. The appendix answering Russell, which opens on how unwelcome it is for a scientific writer to have a foundation shaken after the work is finished. \href{https://mathshistory.st-andrews.ac.uk/Bookpages/Frege10.html}{MacTutor extract}.
\par\medskip

\noindent Peter D. Gr\"unwald (2005). A tutorial introduction to the minimum description length principle. In Peter D. Gr\"unwald, In Jae Myung and Mark A. Pitt (eds.), Advances in Minimum Description Length: Theory and Applications, MIT Press, Cambridge, Massachusetts, pp. 3--80, as the chapters Introducing the Minimum Description Length Principle and Minimum Description Length Tutorial. \href{https://arxiv.org/abs/math/0406077}{arXiv: math/0406077}. The length of a model depends on the code chosen for it, so the code needs a principle behind it.
\par\medskip

\noindent Paul R. Halmos (1960). Naive Set Theory. D. Van Nostrand, Princeton; Springer reprint, 1974. \href{https://doi.org/10.1007/978-1-4757-1645-0}{DOI: 10.1007/978-1-4757-1645-0}. Section 7, on relations: equivalence relations and partitions are the same thing.
\par\medskip

\noindent Geoffrey Hellman (1989). Mathematics without Numbers: Towards a Modal-Structural Interpretation. Clarendon Press, Oxford. \href{https://doi.org/10.1093/0198240341.001.0001}{DOI: 10.1093/0198240341.001.0001}.
\par\medskip

\noindent Carl Huffman (2003, revised 2024). Philolaus. Stanford Encyclopedia of Philosophy. \href{https://plato.stanford.edu/entries/philolaus/}{Entry}. Philolaus's book as probably the first written by a Pythagorean, and the long argument over which of his fragments are genuine.
\par\medskip

\noindent Carl Huffman (2006, revised 2024). Pythagoreanism. Stanford Encyclopedia of Philosophy. \href{https://plato.stanford.edu/entries/pythagoreanism/}{Entry}. Following Burkert, no ancient source connects Hippasus to the discovery of the irrational.
\par\medskip

\noindent Aristotle. Prior Analytics. Translated by A. J. Jenkinson. Book I, chapter 23 (41a23--30): the diagonal is incommensurate with the side, because odd numbers would equal evens if it were commensurate. \href{http://classics.mit.edu/Aristotle/prior.mb.txt}{MIT Classics text}.
\par\medskip

\noindent Aristotle. Metaphysics. Translated by W. D. Ross. Book I, chapter 5 (985b--986a), on the Pythagoreans: ``they supposed the elements of numbers to be the elements of all things'' (\href{http://classics.mit.edu/Aristotle/metaphysics.1.i.html}{MIT Classics, Book I}); Book I, chapter 9 (990b17), where the ``third man'' is named; Book VII, chapter 3, the bronze as matter, the shape as form and the statue as the compound; Book VII, chapter 11 (1036a31 onward), the circle in bronze, stone or wood, and the warning that an animal's parts are parts only when they can do their work (\href{http://classics.mit.edu/Aristotle/metaphysics.7.vii.html}{MIT Classics, Book VII}).
\par\medskip

\noindent Iamblichus (c.\ AD 300). On the Pythagorean Life. Translated by Thomas Taylor as Iamblichus' Life of Pythagoras, A. J. Valpy, London, 1818. \href{https://www.gutenberg.org/ebooks/63300}{Project Gutenberg 63300}. Hippasus drowned for the dodecahedron in chapter XVIII; the unnamed man who let out the incommensurable in chapter XXXIV.
\par\medskip

\noindent Yonggun Jun, Mom\v{c}ilo Gavrilov and John Bechhoefer (2014). High-precision test of Landauer's principle in a feedback trap. Physical Review Letters 113, 190601. \href{https://doi.org/10.1103/PhysRevLett.113.190601}{DOI: 10.1103/PhysRevLett.113.190601}. Erasure at $0.71 \pm 0.03\,kT$ and the control that squashes nothing at $0.05 \pm 0.03\,kT$, on the same bead.
\par\medskip

\noindent Felix Klein (1872). Vergleichende Betrachtungen \"uber neuere geometrische Forschungen (the Erlangen program). Andreas Deichert, Erlangen. English translation by M. W. Haskell, A comparative review of recent researches in geometry, Bulletin of the New York Mathematical Society 2 (1893), pp. 215--249. \href{https://doi.org/10.1090/S0002-9904-1893-00147-X}{DOI: 10.1090/S0002-9904-1893-00147-X}.
\par\medskip

\noindent Ralf Kr\"omer (2006). La ``machine de Grothendieck'' se fonde-t-elle seulement sur des vocables m\'etamath\'ematiques? Bourbaki et les cat\'egories au cours des ann\'ees cinquante. Revue d'histoire des math\'ematiques 12(1), pp. 119--162. \href{https://doi.org/10.24033/rhm.38}{DOI: 10.24033/rhm.38}. Bourbaki's attempt at categories in the 1950s.
\par\medskip

\noindent James Ladyman (1998). What is structural realism? Studies in History and Philosophy of Science Part A 29(3), pp. 409--424. \href{https://doi.org/10.1016/S0039-3681(98)80129-5}{DOI: 10.1016/S0039-3681(98)80129-5}.
\par\medskip

\noindent Saunders Mac Lane (1971; 2nd edition 1998). Categories for the Working Mathematician. Graduate Texts in Mathematics 5, Springer, New York. \href{https://doi.org/10.1007/978-1-4757-4721-8}{DOI: 10.1007/978-1-4757-4721-8}. Where the words came from: category from Aristotle and Kant, functor from Carnap.
\par\medskip

\noindent Gottfried Wilhelm Leibniz (1679). Letter to Christiaan Huygens, 8 September 1679, with the essay on the analysis of situation. In C. I. Gerhardt (ed.), Leibnizens mathematische Schriften, vol. 2, A. Asher, Berlin, 1850, from p.~17; Huygens's reply, which Gerhardt dates to 22 November 1679, from p.~27. \href{https://archive.org/details/leibnizensmathe08leibgoog}{Internet Archive}. English in L. E. Loemker (ed.), Philosophical Papers and Letters, 2nd edition, Reidel, 1969, pp. 248--258, \href{https://doi.org/10.1007/978-94-010-1426-7_28}{DOI: 10.1007/978-94-010-1426-7\_28}.
\par\medskip

\noindent David Lewis (1983). New work for a theory of universals. Australasian Journal of Philosophy 61(4), pp. 343--377. \href{https://doi.org/10.1080/00048408312341131}{DOI: 10.1080/00048408312341131}. The answer to Putnam through natural properties, which Lewis credits to Merrill on p.~371.
\par\medskip

\noindent MacTutor History of Mathematics. Emmy Amalie Noether. University of St Andrews. \href{https://mathshistory.st-andrews.ac.uk/Biographies/Noether_Emmy/}{Biography}. Habilitation closed to women, her courses advertised under Hilbert's name, and habilitation at last in 1919.
\par\medskip

\noindent MacTutor History of Mathematics. Farkas Wolfgang Bolyai. University of St Andrews. \href{https://mathshistory.st-andrews.ac.uk/Biographies/Bolyai_Farkas/}{Biography}. His letter of 4 April 1820 begging his son to leave the parallels alone.
\par\medskip

\noindent Michael Makkai (1998). Towards a categorical foundation of mathematics. In J. A. Makowsky and E. V. Ravve (eds.), Logic Colloquium '95, Lecture Notes in Logic 11, Springer, Berlin, pp. 153--190. \href{https://doi.org/10.1007/978-3-662-22108-2_11}{DOI: 10.1007/978-3-662-22108-2\_11}. ``No global equality present for all things simultaneously'' and a set as ``a relatively orderly part of the world in which an equality predicate reigns,'' both on p.~156.
\par\medskip

\noindent G. H. Merrill (1980). The model-theoretic argument against realism. Philosophy of Science 47(1), pp. 69--81. \href{https://doi.org/10.1086/288910}{DOI: 10.1086/288910}.
\par\medskip

\noindent Augustus De Morgan (1864). On the Syllogism, No. IV, and on the Logic of Relations. Transactions of the Cambridge Philosophical Society 10, pp. 331--358. \href{https://archive.org/details/transactionsofca10camb}{Internet Archive}. Signed at University College London on 12 November 1859, with an appendix dated 15 March 1860, and read on 23 April 1860.
\par\medskip

\noindent Augustus De Morgan (1966). On the Syllogism and Other Logical Writings. Edited by Peter Heath. Routledge and Kegan Paul, London. The fourth paper on the syllogism is on pp. 208--246.
\par\medskip

\noindent Emmy Noether (1918). Invariante Variationsprobleme. Nachrichten von der Gesellschaft der Wissenschaften zu G\"ottingen, Mathematisch-Physikalische Klasse, pp. 235--257. \href{https://eudml.org/doc/59024}{EuDML 59024}.
\par\medskip

\noindent Charles S. Peirce (1870, printed 1873). Description of a notation for the logic of relatives, resulting from an amplification of the conceptions of Boole's calculus of logic. Memoirs of the American Academy of Arts and Sciences 9(2), pp. 317--378. \href{https://doi.org/10.2307/25058006}{DOI: 10.2307/25058006}. Boole's algebra carried over to relations, with examples like ``lover of a servant of'' a woman.
\par\medskip

\noindent Charles S. Peirce (1879). Note on the theory of the economy of research. Report of the Superintendent of the United States Coast Survey for the year ending June 1876, Appendix 14, pp. 197--201. Government Printing Office, Washington. \href{https://archive.org/details/1876-us-coast-survey-report}{Internet Archive}. ``The relations between the utility and the cost of diminishing the probable error of our knowledge.''
\par\medskip

\noindent Plato. Parmenides 132a--b. Translated by H. N. Fowler. \href{https://www.perseus.tufts.edu/hopper/text?doc=Perseus:text:1999.01.0174:text=Parm.:section=132a}{Perseus text}. The regress of greatness that Aristotle later called the third man.
\par\medskip

\noindent Plato. Timaeus 50c--56b. Translated by W. R. M. Lamb. \href{https://www.perseus.tufts.edu/hopper/text?doc=Perseus:text:1999.01.0180}{Perseus text}. The receptacle that takes impressions, the two right triangles, and the four elements as regular solids.
\par\medskip

\noindent Proclus. A Commentary on the First Book of Euclid's Elements. Translated by Glenn R. Morrow. Princeton University Press, 1970. His own attempt to prove the parallel postulate: \href{https://mathshistory.st-andrews.ac.uk/Extras/Proclus_parallel_postulate/}{MacTutor extract}.
\par\medskip

\noindent The Univalent Foundations Program (2013). Homotopy Type Theory: Univalent Foundations of Mathematics. Institute for Advanced Study, Princeton. \href{https://homotopytypetheory.org/book/}{Book}; \href{https://arxiv.org/abs/1308.0729}{arXiv: 1308.0729}. Written during the 2012--13 special year organized by Steve Awodey, Thierry Coquand and Vladimir Voevodsky; univalence as ``identity is equivalent to equivalence'' is in the introduction.
\par\medskip

\noindent Hilary Putnam (1980). Models and reality. Journal of Symbolic Logic 45(3), pp. 464--482. \href{https://doi.org/10.2307/2273415}{DOI: 10.2307/2273415}. The model-theoretic argument Lewis answers.
\par\medskip

\noindent Michael D. Resnik (1997). Mathematics as a Science of Patterns. Clarendon Press, Oxford. \href{https://doi.org/10.1093/0198250142.001.0001}{DOI: 10.1093/0198250142.001.0001}.
\par\medskip

\noindent Jorma Rissanen (1978). Modeling by shortest data description. Automatica 14(5), pp. 465--471. \href{https://doi.org/10.1016/0005-1098(78)90005-5}{DOI: 10.1016/0005-1098(78)90005-5}. The minimum description length principle.
\par\medskip

\noindent Bertrand Russell (1902). Letter to Gottlob Frege, 16 June 1902, and Frege's reply. In Jean van Heijenoort (ed.), From Frege to G\"odel: A Source Book in Mathematical Logic, 1879--1931, Harvard University Press, 1967, pp. 124--125 and 127--128. \href{https://plato.stanford.edu/entries/russell-paradox/}{Stanford Encyclopedia of Philosophy, Russell's paradox}.
\par\medskip

\noindent Girolamo Saccheri (1733). Euclides ab omni naevo vindicatus. Milan. English translation by George Bruce Halsted, Girolamo Saccheri's Euclides Vindicatus, Open Court, Chicago, 1920. \href{https://archive.org/details/girolamosacceri00saccuoft}{Internet Archive}. Proposition XXXIII calls the hypothesis of the acute angle ``repugnant to the nature of the straight line.''
\par\medskip

\noindent Claude E. Shannon (1951). Prediction and entropy of printed English. Bell System Technical Journal 30(1), pp. 50--64. \href{https://doi.org/10.1002/j.1538-7305.1951.tb01366.x}{DOI: 10.1002/j.1538-7305.1951.tb01366.x}. The game of guessing the next letter.
\par\medskip

\noindent Stewart Shapiro (1997). Philosophy of Mathematics: Structure and Ontology. Oxford University Press, New York. \href{https://doi.org/10.1093/0195139305.001.0001}{DOI: 10.1093/0195139305.001.0001}.
\par\medskip

\noindent Theodore Sider (2011). Writing the Book of the World. Oxford University Press, Oxford. \href{https://doi.org/10.1093/acprof:oso/9780199697908.001.0001}{DOI: 10.1093/acprof:oso/9780199697908.001.0001}. ``The world has a distinguished structure, a privileged description'' is in the preface, and ``structure is a worldly, not conceptual or linguistic, matter'' in chapter 1, as read in \href{https://www.tedsider.org/books/wbotw_sample.pdf}{the author's sample}.
\par\medskip

\noindent Stanley Burris and Marcel Jackson. George Boole. Stanford Encyclopedia of Philosophy. \href{https://plato.stanford.edu/entries/boole/}{Entry}. Boole as a schoolmaster who carried on his own education.
\par\medskip

\noindent Alfred Tarski (1954). Contributions to the theory of models, I and II. Indagationes Mathematicae 16, pp. 572--581 and 582--588. \href{https://doi.org/10.1016/S1385-7258(54)50074-0}{DOI: 10.1016/S1385-7258(54)50074-0}; \href{https://doi.org/10.1016/S1385-7258(54)50075-2}{DOI: 10.1016/S1385-7258(54)50075-2}. The first part announces the field under the name theory of models.
\par\medskip

\noindent Alfred Tarski (1986). What are logical notions? Edited by John Corcoran. History and Philosophy of Logic 7(2), pp. 143--154. \href{https://doi.org/10.1080/01445348608837096}{DOI: 10.1080/01445348608837096}. A lecture given at Bedford College, London, on 16 May 1966: logical notions are those invariant under every one-to-one transformation of the world onto itself.
\par\medskip

\noindent Nikolai K. Vereshchagin and Paul M. B. Vit\'anyi (2004). Kolmogorov's structure functions and model selection. IEEE Transactions on Information Theory 50(12), pp. 3265--3290. \href{https://doi.org/10.1109/TIT.2004.838346}{DOI: 10.1109/TIT.2004.838346}; \href{https://arxiv.org/abs/cs/0204037}{arXiv: cs/0204037}. ``Meaningless information, random noise'' is in the section on meaningful information, as read in the arXiv version.
\par\medskip

\noindent Abraham Wald (1945). Sequential tests of statistical hypotheses. Annals of Mathematical Statistics 16(2), pp. 117--186. \href{https://doi.org/10.1214/aoms/1177731118}{DOI: 10.1214/aoms/1177731118}.
\par\medskip

\subsection*{Chapter 11. The counter}

\noindent Lorenz Breidenbach, Phil Daian and Florian Tram\`er, with advice from Ari Juels (2017). GasToken. \href{http://web.archive.org/web/20180308200043/https://gastoken.io/}{Archived site}; \href{https://github.com/projectchicago/gastoken}{code}. Storing gas when it is cheap and spending the refunds when it is dear.
\par\medskip

\noindent Vitalik Buterin and Martin Swende (2021). EIP-3529: Reduction in refunds. Ethereum Improvement Proposals. \href{https://eips.ethereum.org/EIPS/eip-3529}{eips.ethereum.org/EIPS/eip-3529}. Refunds capped at a fifth of the gas used, with GasToken and blocks reaching nearly twice their gas limit given as reasons.
\par\medskip

\noindent Maurice Clint (1973). Program proving: Coroutines. Acta Informatica 2(1), pp. 50--63. \href{https://doi.org/10.1007/BF00571463}{DOI: 10.1007/BF00571463}. The earliest use of auxiliary variables Owicki and Gries had found.
\par\medskip

\noindent Bob Coecke, Tobias Fritz and Robert W. Spekkens (2016). A mathematical theory of resources. Information and Computation 250, pp. 59--86. \href{https://doi.org/10.1016/j.ic.2016.02.008}{DOI: 10.1016/j.ic.2016.02.008}.
\par\medskip

\noindent Thomas Colcombet and Daniela Petri\c{s}an (2020). Automata minimization: a functorial approach. Logical Methods in Computer Science 16(1), paper 32. \href{https://doi.org/10.23638/LMCS-16(1:32)2020}{DOI: 10.23638/LMCS-16(1:32)2020}.
\par\medskip

\noindent Norman Danner, Daniel R. Licata and Ramyaa Ramyaa (2015). Denotational cost semantics for functional languages with inductive types. Proceedings of ICFP 2015, pp. 140--151. \href{https://doi.org/10.1145/2784731.2784749}{DOI: 10.1145/2784731.2784749}.
\par\medskip

\noindent Norman Danner and Daniel R. Licata (2022). Denotational semantics as a foundation for cost recurrence extraction for functional languages. Journal of Functional Programming 32, e8. \href{https://doi.org/10.1017/S095679682200003X}{DOI: 10.1017/S095679682200003X}.
\par\medskip

\noindent Jean-Christophe Filli\^atre, L\'eon Gondelman and Andrei Paskevich (2014). The spirit of ghost code. Computer Aided Verification (CAV 2014), Lecture Notes in Computer Science 8559, Springer, pp. 1--16. \href{https://doi.org/10.1007/978-3-319-08867-9_1}{DOI: 10.1007/978-3-319-08867-9\_1}. When ghost code can be erased with no difference in what the program computes.
\par\medskip

\noindent Tobias Fritz (2017). Resource convertibility and ordered commutative monoids. Mathematical Structures in Computer Science 27(6), pp. 850--938. \href{https://doi.org/10.1017/S0960129515000444}{DOI: 10.1017/S0960129515000444}; \href{https://arxiv.org/abs/1504.03661}{arXiv: 1504.03661}.
\par\medskip

\noindent Joseph A. Goguen (1972). Realization is universal. Mathematical Systems Theory 6(4), pp. 359--374. \href{https://doi.org/10.1007/BF01843493}{DOI: 10.1007/BF01843493}.
\par\medskip

\noindent Tom\'a\v{s} Gonda and Robert W. Spekkens (2023). Monotones in general resource theories. Compositionality 5, article 7. \href{https://doi.org/10.32408/compositionality-5-7}{DOI: 10.32408/compositionality-5-7}.
\par\medskip

\noindent Geoffrey M. Hodgson (n.d.). The secret history of the tally stick. \href{https://www.geoffreymhodgson.uk/secret-history-of-tally-stick}{Web page}. Tallies under Henry I from about 1100 until 1826, the stock and the foil, and the stockholder.
\par\medskip

\noindent Jan Hoffmann, Klaus Aehlig and Martin Hofmann (2011). Multivariate amortized resource analysis. Proceedings of POPL 2011, pp. 357--370. \href{https://doi.org/10.1145/1926385.1926427}{DOI: 10.1145/1926385.1926427}. Journal version: ACM Transactions on Programming Languages and Systems 34(3), 2012, article 14, \href{https://doi.org/10.1145/2362389.2362393}{DOI: 10.1145/2362389.2362393}.
\par\medskip

\noindent Martin Hofmann and Steffen Jost (2003). Static prediction of heap space usage for first-order functional programs. Proceedings of POPL 2003, pp. 185--197. \href{https://doi.org/10.1145/604131.604148}{DOI: 10.1145/604131.604148}. The potential method put into the checking of programs.
\par\medskip

\noindent Paul R. Johnson and Robert H. Thomas (1975). The Maintenance of Duplicate Databases. RFC 677, 27 January 1975. \href{https://www.rfc-editor.org/rfc/rfc677.html}{RFC 677}. The timestamps Lamport started from.
\par\medskip

\noindent Shin-ya Katsumata (2014). Parametric effect monads and semantics of effect systems. Proceedings of POPL 2014, pp. 633--645. \href{https://doi.org/10.1145/2535838.2535846}{DOI: 10.1145/2535838.2535846}.
\par\medskip

\noindent Narayana R. Kocherlakota (1998). Money is memory. Journal of Economic Theory 81(2), pp. 232--251. \href{https://doi.org/10.1006/jeth.1997.2357}{DOI: 10.1006/jeth.1997.2357}. Where nobody can commit, anything money makes possible, memory of past dealings makes possible too.
\par\medskip

\noindent Leslie Lamport (1978). Time, clocks, and the ordering of events in a distributed system. Communications of the ACM 21(7), pp. 558--565. \href{https://doi.org/10.1145/359545.359563}{DOI: 10.1145/359545.359563}. ``Counters with no actual timing mechanism'' is on p.~559.
\par\medskip

\noindent Leslie Lamport (n.d.). The Writings of Leslie Lamport, note on ``Time, Clocks and the Ordering of Events in a Distributed System.'' \href{https://lamport.azurewebsites.net/pubs/pubs.html}{Web page}. The Johnson and Thomas note, special relativity, and Jim Gray's two opinions of the paper.
\par\medskip

\noindent Elliott H. Lieb and Jakob Yngvason (1999). The physics and mathematics of the second law of thermodynamics. Physics Reports 310(1), pp. 1--96. \href{https://doi.org/10.1016/S0370-1573(98)00082-9}{DOI: 10.1016/S0370-1573(98)00082-9}. From an order of which states can reach which, an entropy that is unique once zero and unit are chosen.
\par\medskip

\noindent Aleksandr M. Lyapunov (1892). The general problem of the stability of motion. Doctoral thesis, Kharkov (in Russian). French translation by \'Edouard Davaux, naval engineer at Toulon, revised and corrected by the author: Probl\`eme g\'en\'eral de la stabilit\'e du mouvement, Annales de la Facult\'e des Sciences de Toulouse, 2e s\'erie, 9 (1907), pp. 203--474, \href{http://www.numdam.org/item/AFST_1907_2_9__203_0/}{Numdam}. English translation by A. T. Fuller, International Journal of Control 55(3) (1992), pp. 531--772, \href{https://doi.org/10.1080/00207179208934253}{DOI: 10.1080/00207179208934253} (opening part).
\par\medskip

\noindent Alfred Mitchell-Innes (1913). What is money? The Banking Law Journal, May 1913, pp. 377--408. \href{https://www.levyinstitute.org/pubs/wp_717.pdf}{Levy Institute working paper discussing it}. The crown paying its creditors by ``raising a tally'' on some branch of revenue, p.~398.
\par\medskip

\noindent Richard fitz Nigel (c.\ 1180). The Dialogue Concerning the Exchequer (Dialogus de Scaccario). Translated in Ernest F. Henderson, Select Historical Documents of the Middle Ages, George Bell and Sons, London, 1896. \href{https://avalon.law.yale.edu/medieval/excheq.asp}{Avalon Project text}. Modern edition and translation by Emilie Amt, Dialogus de Scaccario, and Constitutio Domus Regis, Oxford Medieval Texts, Oxford University Press, 2007. The widths of the notches on a tally, from a palm for a thousand pounds down to a bare cut for a penny.
\par\medskip

\noindent Yue Niu, Jonathan Sterling, Harrison Grodin and Robert Harper (2022). A cost-aware logical framework. Proceedings of the ACM on Programming Languages 6(POPL), article 9. \href{https://doi.org/10.1145/3498670}{DOI: 10.1145/3498670}. The framework calf.
\par\medskip

\noindent Dominic Orchard, Vilem-Benjamin Liepelt and Harley Eades III (2019). Quantitative program reasoning with graded modal types. Proceedings of the ACM on Programming Languages 3(ICFP), article 110. \href{https://doi.org/10.1145/3341714}{DOI: 10.1145/3341714}.
\par\medskip

\noindent Susan Owicki and David Gries (1976). An axiomatic proof technique for parallel programs I. Acta Informatica 6(4), pp. 319--340. \href{https://doi.org/10.1007/BF00268134}{DOI: 10.1007/BF00268134}. Auxiliary variables, needed for the proof and not for the program, which only ever get assigned.
\par\medskip

\noindent Marcel-Paul Sch\"utzenberger (1961). On the definition of a family of automata. Information and Control 4(2--3), pp. 245--270. \href{https://doi.org/10.1016/S0019-9958(61)80020-X}{DOI: 10.1016/S0019-9958(61)80020-X}. Automata with weights on their moves.
\par\medskip

\noindent Caroline Shenton (2012). The Day Parliament Burned Down. Oxford University Press, Oxford. The tallies burned in the furnaces under the House of Lords on 16 October 1834, and the fire that started in the flues.
\par\medskip

\noindent Robert Endre Tarjan (1985). Amortized computational complexity. SIAM Journal on Algebraic and Discrete Methods 6(2), pp. 306--318. \href{https://doi.org/10.1137/0606031}{DOI: 10.1137/0606031}. The potential method, with the physicist's view credited to Daniel Sleator, and a stack whose total work is at most twice the number of pushes.
\par\medskip

\noindent Alan M. Turing (1949). Checking a large routine. In Report of a Conference on High Speed Automatic Calculating Machines, University Mathematical Laboratory, Cambridge, pp. 67--69. \href{https://turingarchive.kings.cam.ac.uk/publications-lectures-and-talks-amtb/amt-b-8}{Turing Digital Archive}. Reprinted by F. L. Morris and C. B. Jones, An early program proof by Alan Turing, Annals of the History of Computing 6(2), 1984, pp. 139--143, \href{https://doi.org/10.1109/MAHC.1984.10017}{DOI: 10.1109/MAHC.1984.10017}. A quantity that decreases continually and vanishes when the machine stops.
\par\medskip

\noindent UK Parliament. The Great Fire of 1834. Living Heritage. \href{https://www.parliament.uk/about/living-heritage/building/palace/architecture/palacestructure/great-fire/}{Web page}. Two cart-loads of tally sticks, unused since 1826, burned in the stoves under the House of Lords.
\par\medskip

\noindent Gavin Wood (2014). Ethereum: A Secure Decentralised Generalised Transaction Ledger (the yellow paper). \href{https://ethereum.github.io/yellowpaper/paper.pdf}{Current version}. Why every computation costs a fee, and the gas table in its appendix.
\par\medskip

\subsection*{Chapter 12. What a ticked box lets you lean on}

\noindent Sergei N. Artemov (2001). Explicit provability and constructive semantics. Bulletin of Symbolic Logic 7(1), pp. 1--36. \href{https://doi.org/10.2307/2687821}{DOI: 10.2307/2687821}. The Logic of Proofs, where knowing comes with the reason written out.
\par\medskip

\noindent A. J. Ayer (1956). The Problem of Knowledge. Macmillan, London, p.~34. One of the two statements of the definition that Gettier takes apart.
\par\medskip

\noindent Alexandru Baltag, Bryan Renne and Sonja Smets (2014). The logic of justified belief, explicit knowledge, and conclusive evidence. Annals of Pure and Applied Logic 165(1), pp. 49--81. \href{https://doi.org/10.1016/j.apal.2013.07.005}{DOI: 10.1016/j.apal.2013.07.005}.
\par\medskip

\noindent Walter J. Boyne (2013). The Checklist. Air Force Magazine, August 2013, pp. 52--56. \href{https://www.airandspaceforces.com/article/0813checklist/}{Article}. The Model 299 crash of 30 October 1935, the written pilot routines going back to 1911, and the judgement that ``it took the Model 299 crash to institutionalize its use.''
\par\medskip

\noindent Robert B. Brandom (1994). Making It Explicit: Reasoning, Representing, and Discursive Commitment. Harvard University Press, Cambridge, Massachusetts. Commitment, entitlement, and the score speakers keep on each other; his home page at the University of Pittsburgh is \href{https://sites.pitt.edu/~rbrandom/}{sites.pitt.edu/\textasciitilde rbrandom}.
\par\medskip

\noindent Martin Caidin (1968). Flying Forts. Meredith Press, New York. Has ``rumors that the 299 was too much airplane for any one pilot.''
\par\medskip

\noindent Roderick M. Chisholm (1957). Perceiving: A Philosophical Study. Cornell University Press, Ithaca, p.~16. One of the two statements of the definition that Gettier takes apart.
\par\medskip

\noindent DeWitt S. Copp (1980). A Few Great Captains: The Men and Events That Shaped the Development of U.S. Air Power. Doubleday, Garden City, New York. The earliest place found with the exact line ``too much airplane for one man to fly,'' given as word from the War Department.
\par\medskip

\noindent Charles Dickens (1843). A Christmas Carol. In Prose. Being a Ghost Story of Christmas. Chapman and Hall, London. Stave Three, the game of Yes and No at Fred's party. \href{https://www.gutenberg.org/ebooks/46}{Project Gutenberg text}.
\par\medskip

\noindent Julien Dutant (2015). The Legend of the Justified True Belief Analysis. Philosophical Perspectives 29(1), pp. 95--145. \href{https://doi.org/10.1111/phpe.12061}{DOI: 10.1111/phpe.12061}. Known here only from its abstract and summaries.
\par\medskip

\noindent Euclid. Elements, Book X. \href{https://mathcs.clarku.edu/~djoyce/elements/bookX/bookX.html}{D. E. Joyce's online edition}. The book on lengths that share no common measure, thought to be largely Theaetetus's work.
\par\medskip

\noindent Plato. Theaetetus. Trans. Harold N. Fowler, Plato in Twelve Volumes, vol. 12, Harvard University Press, 1921. \href{http://www.perseus.tufts.edu/hopper/text?doc=Perseus:text:1999.01.0172}{Perseus text}. Theaetetus carried home from the army at Corinth and Euclides's notes corrected by Socrates (142a--143c), Theodorus's square roots up to seventeen (147d--148b), and the three meanings of an account (201c--210b).
\par\medskip

\noindent Peter G\"ardenfors (1986). Belief revisions and the Ramsey test for conditionals. The Philosophical Review 95(1), pp. 81--93. \href{https://doi.org/10.2307/2185133}{DOI: 10.2307/2185133}. The impossibility theorem.
\par\medskip

\noindent Atul Gawande (2007). The Checklist. The New Yorker, 10 December 2007. \href{https://www.newyorker.com/magazine/2007/12/10/the-checklist}{Article}. Gives the line ``too much airplane for one man to fly'' as a newspaper's, without naming the paper.
\par\medskip

\noindent G\"osta Grahne (1998). Updates and counterfactuals. Journal of Logic and Computation 8(1), pp. 87--117. \href{https://doi.org/10.1093/logcom/8.1.87}{DOI: 10.1093/logcom/8.1.87}. G\"ardenfors's theorem bites revision, not update.
\par\medskip

\noindent Ralph V. L. Hartley (1928). Transmission of Information. Bell System Technical Journal 7(3), pp. 535--563. \href{https://doi.org/10.1002/j.1538-7305.1928.tb01236.x}{DOI: 10.1002/j.1538-7305.1928.tb01236.x}. Information measured by the count of possibilities alone.
\par\medskip

\noindent David Hume (1748). An Enquiry Concerning Human Understanding, Section X, ``Of Miracles,'' paragraph 4: ``A wise man, therefore, proportions his belief to the evidence.'' \href{https://davidhume.org/texts/e/10}{Hume Texts Online, E 10.4}.
\par\medskip

\noindent Jonathan Jenkins Ichikawa and Matthias Steup (2001, revised 2026). The Analysis of Knowledge. Stanford Encyclopedia of Philosophy. \href{https://plato.stanford.edu/entries/knowledge-analysis/}{Entry}. The ``traditional'' analysis and the ``convenient fiction'' that it was widely held, and the older cases from Dharmottara and Peter of Mantua.
\par\medskip

\noindent Edward Jablonski (1965). Flying Fortress: The Illustrated Biography of the B-17s and the Men Who Flew Them. Doubleday, Garden City, New York. The talk after the crash that the airplane was ``too much airplane for one pilot to handle.''
\par\medskip

\noindent Lester R. Ford Jr. and Selmer M. Johnson (1959). A Tournament Problem. American Mathematical Monthly 66(5), pp. 387--389. \href{https://doi.org/10.1080/00029890.1959.11989306}{DOI: 10.1080/00029890.1959.11989306}. Merge insertion.
\par\medskip

\noindent Hirofumi Katsuno and Alberto O. Mendelzon (1991). On the difference between updating a knowledge base and revising it. Proceedings of the Second International Conference on Principles of Knowledge Representation and Reasoning (KR'91), pp. 387--394; reprinted in Peter G\"ardenfors (ed.), Belief Revision, Cambridge University Press (1992), pp. 183--203. \href{https://doi.org/10.1017/CBO9780511526664.007}{DOI: 10.1017/CBO9780511526664.007}.
\par\medskip

\noindent Arthur M. Keller and Marianne Winslett Wilkins (1985). On the use of an extended relational model to handle changing incomplete information. IEEE Transactions on Software Engineering SE-11(7), pp. 620--633. \href{https://doi.org/10.1109/TSE.1985.232506}{DOI: 10.1109/TSE.1985.232506}.
\par\medskip

\noindent Donald E. Knuth (1998). The Art of Computer Programming, Volume 3: Sorting and Searching, second edition. Addison-Wesley. Section 5.3.1: the counting floor for sorting by comparisons, merge insertion, and Wells's result for twelve items.
\par\medskip

\noindent Isaac Levi (1980). The Enterprise of Knowledge. MIT Press, Cambridge, Massachusetts. Where routine and deliberate expansion first appear.
\par\medskip

\noindent Isaac Levi (1991). The Fixation of Belief and Its Undoing: Changing Beliefs through Inquiry. Cambridge University Press. \href{https://doi.org/10.1017/CBO9780511663819}{DOI: 10.1017/CBO9780511663819}. Routine and deliberate expansion.
\par\medskip

\noindent George C. Necula and Peter Lee (1996). Safe Kernel Extensions Without Run-Time Checking. Proceedings of the 2nd USENIX Symposium on Operating Systems Design and Implementation (OSDI '96), pp. 229--243. \href{https://www.usenix.org/conference/osdi-96/safe-kernel-extensions-without-run-time-checking}{USENIX page}.
\par\medskip

\noindent George C. Necula (1997). Proof-Carrying Code. Proceedings of the 24th ACM SIGPLAN-SIGACT Symposium on Principles of Programming Languages (POPL '97), pp. 106--119. \href{https://doi.org/10.1145/263699.263712}{DOI: 10.1145/263699.263712}.
\par\medskip

\noindent New Castle News (1935). Report on the crash of the Boeing bomber at Wright Field, 12 November 1935, New Castle, Pennsylvania. One of the papers that November reporting the locked controls.
\par\medskip

\noindent J. J. O'Connor and E. F. Robertson (1999). Theaetetus of Athens. MacTutor History of Mathematics archive, University of St Andrews. \href{https://mathshistory.st-andrews.ac.uk/Biographies/Theaetetus/}{Biography}. His death in about 369 BC after being wounded in the fighting at Corinth, and Book X of Euclid as largely his work.
\par\medskip

\noindent Bertrand Russell (1948). Human Knowledge: Its Scope and Limits. George Allen and Unwin, London. The stopped clock.
\par\medskip

\noindent Krister Segerberg (1998). Irrevocable Belief Revision in Dynamic Doxastic Logic. Notre Dame Journal of Formal Logic 39(3), pp. 287--306. \href{https://doi.org/10.1305/ndjfl/1039182247}{DOI: 10.1305/ndjfl/1039182247}.
\par\medskip

\noindent Neil Tennant (2008). Belief-revision, the Ramsey test, monotonicity, and the so-called impossibility results. The Review of Symbolic Logic 1(4), pp. 402--423. \href{https://doi.org/10.1017/S1755020308090023}{DOI: 10.1017/S1755020308090023}.
\par\medskip

\noindent Mark B. Wells (1965). Applications of a language for computing in combinatorics. Information Processing 65, Proceedings of IFIP Congress 1965, pp. 497--498. The computer search showing twelve items need thirty comparisons.
\par\medskip

\noindent Linda Zagzebski (1994). The inescapability of Gettier problems. The Philosophical Quarterly 44(174), pp. 65--73. \href{https://doi.org/10.2307/2220147}{DOI: 10.2307/2220147}.
\par\medskip

\subsection*{Chapter 13. What insight is}

\noindent Karl Duncker (1935; English translation 1945). Zur Psychologie des produktiven Denkens. Julius Springer, Berlin. English translation by Lynne S. Lees: On Problem-Solving, Psychological Monographs 58(5), whole no. 270, 1945. \href{https://doi.org/10.1037/h0093599}{DOI: 10.1037/h0093599}. The candle problem.
\par\medskip

\noindent Galileo Galilei (1586). La Bilancetta (The Little Balance). Described at the Galileo Project, Rice University: \href{https://galileo.library.rice.edu/sci/instruments/balance.html}{page}. The words quoted are from the English translation in Laura Fermi and Gilberto Bernardini, Galileo and the Scientific Revolution, Basic Books, 1961, as reproduced at \href{https://math.nyu.edu/Archimedes/Crown/bilancetta.html}{math.nyu.edu}.
\par\medskip

\noindent Sam Glucksberg (1962). The influence of strength of drive on functional fixedness and perceptual recognition. Journal of Experimental Psychology 63(1), pp. 36--41. \href{https://doi.org/10.1037/h0044683}{DOI: 10.1037/h0044683}. The candle problem with prize money on it.
\par\medskip

\noindent I. J. Good (1979). Studies in the History of Probability and Statistics XXXVII: A. M. Turing's statistical work in World War II. Biometrika 66(2), pp. 393--396. \href{https://doi.org/10.1093/biomet/66.2.393}{DOI: 10.1093/biomet/66.2.393}. Turing's expected weight of evidence, the ban and the deciban, and Peirce's 1878 name for the measure.
\par\medskip

\noindent Wolfgang K\"ohler (1917). Intelligenzpr\"ufungen an Anthropoiden. Abhandlungen der K\"oniglich Preussischen Akademie der Wissenschaften, Physikalisch-mathematische Klasse, 1917. English translation by Ella Winter, The Mentality of Apes, Kegan Paul, Trench, Trubner, London, 1925. Sultan and the two sticks.
\par\medskip

\noindent Dennis V. Lindley (1956). On a Measure of the Information Provided by an Experiment. Annals of Mathematical Statistics 27(4), pp. 986--1005. \href{https://doi.org/10.1214/aoms/1177728069}{DOI: 10.1214/aoms/1177728069}.
\par\medskip

\noindent Janet Metcalfe (1986). Premonitions of insight predict impending error. Journal of Experimental Psychology: Learning, Memory, and Cognition 12(4), pp. 623--634. \href{https://doi.org/10.1037/0278-7393.12.4.623}{DOI: 10.1037/0278-7393.12.4.623}. High warmth before an insight answer goes with the answer being wrong.
\par\medskip

\noindent Janet Metcalfe and David Wiebe (1987). Intuition in insight and noninsight problem solving. Memory and Cognition 15(3), pp. 238--246. \href{https://doi.org/10.3758/BF03197722}{DOI: 10.3758/BF03197722}. Warmth that climbs on algebra and stays flat on insight problems.
\par\medskip

\noindent Charles S. Peirce (1878). The Probability of Induction. Popular Science Monthly 12, pp. 705--718. Illustrations of the Logic of Science, fourth paper. \href{https://en.wikisource.org/wiki/Popular_Science_Monthly/Volume_12/April_1878/Illustrations_of_the_Logic_of_Science_IV}{Wikisource text}. Belief measured by a logarithm, so that independent weights of evidence add.
\par\medskip

\noindent Vitruvius. The Ten Books on Architecture, Book IX, Introduction. Trans. Morris Hicky Morgan, Harvard University Press, 1914. \href{https://math.nyu.edu/Archimedes/Crown/Vitruvius.html}{Passage}. Archimedes, Hiero's crown and the bath.
\par\medskip

\subsection*{Chapter 14. A box over any machine}

\noindent Mart\'in Abadi and Leslie Lamport (1991). The Existence of Refinement Mappings. Theoretical Computer Science 82(2), pp. 253--284. \href{https://doi.org/10.1016/0304-3975(91)90224-P}{DOI: 10.1016/0304-3975(91)90224-P}. History variables.
\par\medskip

\noindent David Z. Albert (2000). Time and Chance. Harvard University Press, Cambridge, Massachusetts. The name ``past hypothesis'' for an assumed low-entropy beginning; its roots in Boltzmann are traced in the Stanford Encyclopedia of Philosophy entry Thermodynamic Asymmetry in Time, \href{https://plato.stanford.edu/entries/time-thermo/}{Entry}.
\par\medskip

\noindent Pietro Baroni and Massimiliano Giacomin (2007). On principle-based evaluation of extension-based argumentation semantics. Artificial Intelligence 171(10--15), pp. 675--700. \href{https://doi.org/10.1016/j.artint.2007.04.004}{DOI: 10.1016/j.artint.2007.04.004}.
\par\medskip

\noindent John S. Bell (1975). On wave packet reduction in the Coleman-Hepp model. Helvetica Physica Acta 48, p.~93. Reprinted in Speakable and Unspeakable in Quantum Mechanics, Cambridge University Press, \href{https://www.cambridge.org/core/books/speakable-and-unspeakable-in-quantum-mechanics/on-wave-packet-reduction-in-the-colemanhepp-model/95A04BE1845F5490A462D929930CA7E8}{chapter page}. The reduction happens only in a limit that never comes.
\par\medskip

\noindent Thomas Blanchard (2024). Causation and the Time-Asymmetry of Knowledge. Australasian Journal of Philosophy 102(4), pp. 959--977. \href{https://doi.org/10.1080/00048402.2024.2328629}{DOI: 10.1080/00048402.2024.2328629}.
\par\medskip

\noindent John P. Burgess (1986). The truth is never simple. Journal of Symbolic Logic 51(3), pp. 663--681. \href{https://doi.org/10.2307/2274021}{DOI: 10.2307/2274021}. Quasi-inductive definitions.
\par\medskip

\noindent Martin Caminada (2014). Strong admissibility revisited. In Simon Parsons, Nir Oren, Chris Reed and Federico Cerutti (eds.), Computational Models of Argument: Proceedings of COMMA 2014, Frontiers in Artificial Intelligence and Applications 266, IOS Press, pp. 197--208. \href{https://doi.org/10.3233/978-1-61499-436-7-197}{DOI: 10.3233/978-1-61499-436-7-197}.
\par\medskip

\noindent Wen Tsing Chow and William H. Henrich (1962). Storage matrix. US patent 3,028,659, filed 27 December 1957, granted 10 April 1962. \href{https://patents.google.com/patent/US3028659A/en}{Patent text}. The memory of blown fuses that chapter fourteen describes.
\par\medskip

\noindent Vincent Danos and Jean Krivine (2004). Reversible Communicating Systems. CONCUR 2004, Lecture Notes in Computer Science 3170, pp. 292--307. \href{https://doi.org/10.1007/978-3-540-28644-8_19}{DOI: 10.1007/978-3-540-28644-8\_19}.
\par\medskip

\noindent Vincent Danos and Jean Krivine (2005). Transactions in RCCS. CONCUR 2005, Lecture Notes in Computer Science 3653, pp. 398--412. \href{https://doi.org/10.1007/11539452_31}{DOI: 10.1007/11539452\_31}.
\par\medskip

\noindent David Deutsch and Chiara Marletto (2025). Constructor theory of time. Preprint. \href{https://arxiv.org/abs/2505.08692}{arXiv: 2505.08692}.
\par\medskip

\noindent Rod Downey, Carl Jockusch, Timothy H. McNicholl and Paul Schupp (2015). Asymptotic density and the Ershov hierarchy. Mathematical Logic Quarterly 61(3), pp. 189--195. \href{https://doi.org/10.1002/malq.201300081}{DOI: 10.1002/malq.201300081}; \href{https://arxiv.org/abs/1309.0137}{arXiv: 1309.0137} (revised 2014). Their account of Ershov's hierarchy, where one change of guess from no to yes gives the sets a program can list.
\par\medskip

\noindent Phan Minh Dung (1995). On the acceptability of arguments and its fundamental role in nonmonotonic reasoning, logic programming and n-person games. Artificial Intelligence 77(2), pp. 321--357. \href{https://doi.org/10.1016/0004-3702(94)00041-X}{DOI: 10.1016/0004-3702(94)00041-X}.
\par\medskip

\noindent Aram Ebtekar and Marcus Hutter (2024). Modeling the Arrows of Time with Causal Multibaker Maps. Entropy 26(9), article 776. \href{https://doi.org/10.3390/e26090776}{DOI: 10.3390/e26090776}.
\par\medskip

\noindent W. H. Eccles and F. W. Jordan (1919). A trigger relay utilizing three-electrode thermionic vacuum tubes. The Electrician 83, p.~298, 19 September 1919; reprinted in Radio Review 1(3), pp. 143--146.
\par\medskip

\noindent W. H. Eccles and F. W. Jordan (1920). Improvements in ionic relays. British patent GB 148,582, filed 21 June 1918, published 5 August 1920. \href{https://patents.google.com/patent/GB148582A/en}{Patent record}. The two-valve circuit that became the flip-flop.
\par\medskip

\noindent Yu. L. Ershov (1968). A hierarchy of sets. I. Algebra and Logic 7, pp. 25--43; translated from Algebra i Logika 7(1), pp. 47--74. \href{https://doi.org/10.1007/BF02218750}{DOI: 10.1007/BF02218750}.
\par\medskip

\noindent Murray Gell-Mann (1964). A schematic model of baryons and mesons. Physics Letters 8(3), pp. 214--215. \href{https://doi.org/10.1016/S0031-9163(64)92001-3}{DOI: 10.1016/S0031-9163(64)92001-3}. The paper that named the quark.
\par\medskip

\noindent Murray Gell-Mann and James B. Hartle (1995, printed 1998). Strong decoherence. In D.-H. Feng and B.-L. Hu (eds.), Proceedings of the 4th Drexel Conference on Quantum Non-Integrability: The Quantum-Classical Correspondence, International Press, Hong Kong. \href{https://arxiv.org/abs/gr-qc/9509054}{arXiv: gr-qc/9509054}. The permanence of the past.
\par\medskip

\noindent E. Mark Gold (1965). Limiting recursion. Journal of Symbolic Logic 30(1), pp. 28--48. \href{https://doi.org/10.2307/2270580}{DOI: 10.2307/2270580}.
\par\medskip

\noindent J. J. Halliwell (1999). Arrival Times in Quantum Theory from an Irreversible Detector Model. Progress of Theoretical Physics 102(4), pp. 707--717. \href{https://doi.org/10.1143/PTP.102.707}{DOI: 10.1143/PTP.102.707}. A two-state detector made ``effectively irreversible'' by a large environment.
\par\medskip

\noindent C. L. Hamblin (1962). Translation to and from Polish notation. The Computer Journal 5(3), pp. 210--213. \href{https://doi.org/10.1093/comjnl/5.3.210}{DOI: 10.1093/comjnl/5.3.210}. His reverse Polish notation for computers.
\par\medskip

\noindent C. L. Hamblin (1970). Fallacies. Methuen, London. Chapter 8, formal dialectic, where each speaker has a store of commitments that the rules update.
\par\medskip

\noindent Joel David Hamkins and Andy Lewis (2000). Infinite time Turing machines. Journal of Symbolic Logic 65(2), pp. 567--604. \href{https://doi.org/10.2307/2586556}{DOI: 10.2307/2586556}; \href{https://arxiv.org/abs/math/9808093}{arXiv: math/9808093}.
\par\medskip

\noindent Klaus Havelund and Grigore Ro\c{s}u (2002). Synthesizing monitors for safety properties. Tools and Algorithms for the Construction and Analysis of Systems (TACAS 2002), Lecture Notes in Computer Science 2280, pp. 342--356. \href{https://doi.org/10.1007/3-540-46002-0_24}{DOI: 10.1007/3-540-46002-0\_24}.
\par\medskip

\noindent Klaus Hepp (1972). Quantum theory of measurement and macroscopic observables. Helvetica Physica Acta 45, p.~237. The Coleman-Hepp model.
\par\medskip

\noindent Erik C. W. Krabbe (2001). The problem of retraction in critical discussion. Synthese 127(1--2), pp. 141--159. \href{https://doi.org/10.1023/A:1010318403544}{DOI: 10.1023/A:1010318403544}.
\par\medskip

\noindent Saul Kripke (1975). Outline of a theory of truth. Journal of Philosophy 72(19), pp. 690--716. \href{https://doi.org/10.2307/2024634}{DOI: 10.2307/2024634}. Grounded and ungrounded sentences.
\par\medskip

\noindent Orna Lichtenstein, Amir Pnueli and Lenore Zuck (1985). The glory of the past. Logics of Programs, Lecture Notes in Computer Science 193, pp. 196--218. \href{https://doi.org/10.1007/3-540-15648-8_16}{DOI: 10.1007/3-540-15648-8\_16}. A logic of the past, with the operator ``once.''
\par\medskip

\noindent Robin Milner (1980). A Calculus of Communicating Systems. Lecture Notes in Computer Science 92, Springer. \href{https://doi.org/10.1007/3-540-10235-3}{DOI: 10.1007/3-540-10235-3}.
\par\medskip

\noindent Henry Prakken (2005). Coherence and flexibility in dialogue games for argumentation. Journal of Logic and Computation 15(6), pp. 1009--1040. \href{https://doi.org/10.1093/logcom/exi046}{DOI: 10.1093/logcom/exi046}.
\par\medskip

\noindent Hilary Putnam (1965). Trial and error predicates and the solution to a problem of Mostowski. Journal of Symbolic Logic 30(1), pp. 49--57. \href{https://doi.org/10.2307/2270581}{DOI: 10.2307/2270581}.
\par\medskip

\noindent Wayne Richter and Peter Aczel (1974). Inductive definitions and reflecting properties of admissible ordinals. In J. E. Fenstad and P. G. Hinman (eds.), Generalized Recursion Theory, North-Holland, Studies in Logic and the Foundations of Mathematics 79, pp. 301--381. \href{https://doi.org/10.1016/S0049-237X(08)70592-5}{DOI: 10.1016/S0049-237X(08)70592-5}.
\par\medskip

\noindent Carlo Rovelli (2022). Memory and Entropy. Entropy 24(8), article 1022. \href{https://doi.org/10.3390/e24081022}{DOI: 10.3390/e24081022}; \href{https://arxiv.org/abs/2003.06687}{arXiv: 2003.06687}.
\par\medskip

\noindent Alin Tomescu, Vivek Bhupatiraju, Dimitrios Papadopoulos, Charalampos Papamanthou, Nikos Triandopoulos and Srinivas Devadas (2019). Transparency logs via append-only authenticated dictionaries. Proceedings of the 2019 ACM SIGSAC Conference on Computer and Communications Security (CCS '19), pp. 1299--1316. \href{https://doi.org/10.1145/3319535.3345652}{DOI: 10.1145/3319535.3345652}.
\par\medskip

\noindent Hao Wang (1957). A variant to Turing's theory of computing machines. Journal of the ACM 4(1), pp. 63--92. \href{https://doi.org/10.1145/320856.320867}{DOI: 10.1145/320856.320867}. Machines that mark a square and never erase one.
\par\medskip

\noindent Wikipedia. Programmable ROM, History section. \href{https://en.wikipedia.org/wiki/Programmable_ROM}{Entry}. Chow's invention at Arma in 1956, at the Air Force's request, to hold the targeting constants of the Atlas computer, and ``burn'' as the patent's word.
\par\medskip

\noindent David H. Wolpert (1992). Memory systems, computation, and the second law of thermodynamics. International Journal of Theoretical Physics 31(4), pp. 743--785. \href{https://doi.org/10.1007/BF00670829}{DOI: 10.1007/BF00670829}.
\par\medskip

\noindent David H. Wolpert and Jens Kipper (2024). Memory Systems, the Epistemic Arrow of Time, and the Second Law. Entropy 26(2), article 170. \href{https://doi.org/10.3390/e26020170}{DOI: 10.3390/e26020170}.
\par\medskip

\subsection*{Chapter 16. Reading the box}

\noindent A. P. Ershov (1958). On programming of arithmetic operations. Communications of the ACM 1(8), pp. 3--6. \href{https://doi.org/10.1145/368892.368907}{DOI: 10.1145/368892.368907}. Translated by Morris D. Friedman from Doklady Akademii Nauk SSSR 118(3), 1958, pp. 427--430. Never working out the same expression twice, in time that grows with the number of instructions instead of its square.
\par\medskip

\noindent Kurt G\"odel (1931). \"Uber formal unentscheidbare S\"atze der Principia Mathematica und verwandter Systeme I. Monatshefte f\"ur Mathematik und Physik 38, pp. 173--198. \href{https://doi.org/10.1007/BF01700692}{DOI: 10.1007/BF01700692}.
\par\medskip

\noindent Eiichi Goto (1974). Monocopy and associative algorithms in an extended Lisp. Technical report, Information Science Laboratory, Faculty of Science, University of Tokyo, May 1974. \href{https://www.cs.utexas.edu/~hunt/research/hash-cons/hash-cons-papers/monocopy-goto.pdf}{Scanned copy}. No two copies of the same structure, building on a scheme proposed at the 1973 Programming Symposium.
\par\medskip

\noindent IBM (1974, revised 1975). IBM System/370 Principles of Operation. GA22-7000-4, fourth edition, September 1974, as revised 1 September 1975 by TNL GN22-0498. Programming note to Compare Double and Swap, p.~125. \href{https://ftp.mirrorservice.org/sites/www.bitsavers.org/pdf/ibm/370/princOps/GA22-7000-4_370_Principles_Of_Operation_Sep75.pdf}{Scanned manual}. The list changed behind a program's back, and the change number that brings the risk ``to a negligible level.''
\par\medskip

\noindent Maged M. Michael and Michael L. Scott (1995). Correction of a Memory Management Method for Lock-Free Data Structures. Technical Report 599, Computer Science Department, University of Rochester, December 1995. \href{https://www.cs.rochester.edu/u/scott/papers/1995_TR599.pdf}{Report}. Uses the name ``ABA problem.''
\par\medskip

\noindent Maged M. Michael and Michael L. Scott (1995). Simple, Fast, and Practical Non-Blocking and Blocking Concurrent Queue Algorithms. Technical Report 600, Computer Science Department, University of Rochester, December 1995. \href{https://doi.org/10.21236/ADA309412}{DOI: 10.21236/ADA309412}.
\par\medskip

\noindent Maged M. Michael and Michael L. Scott (1996). Simple, fast, and practical non-blocking and blocking concurrent queue algorithms. Proceedings of the Fifteenth Annual ACM Symposium on Principles of Distributed Computing (PODC '96), pp. 267--275. \href{https://doi.org/10.1145/248052.248106}{DOI: 10.1145/248052.248106}. A value that goes from A to B and back to A, and the counter that catches it.
\par\medskip

\noindent Donald Michie (1968). ``Memo'' Functions and Machine Learning. Nature 218, 6 April 1968, pp. 19--22. \href{https://doi.org/10.1038/218019a0}{DOI: 10.1038/218019a0}. The rule part and the rote part, and each evaluation by rule adding a fresh entry to the rote.
\par\medskip

\noindent R. Kent Treiber (1986). Systems Programming: Coping with Parallelism. Research Report RJ 5118, IBM Almaden Research Center, April 1986. The lock-free stack and its counter beside the head of the list.
\par\medskip

\subsection*{Chapter 17. What a machine can't decide about itself}

\noindent Manuel Blum (1967). A machine-independent theory of the complexity of recursive functions. Journal of the ACM 14(2), pp. 322--336. \href{https://doi.org/10.1145/321386.321395}{DOI: 10.1145/321386.321395}. The speedup theorem: functions with no fastest program.
\par\medskip

\noindent Georg Cantor (1874). \"Uber eine Eigenschaft des Inbegriffes aller reellen algebraischen Zahlen. Journal f\"ur die reine und angewandte Mathematik 77, pp. 258--262. \href{https://doi.org/10.1515/crll.1874.77.258}{DOI: 10.1515/crll.1874.77.258}. The first, harder proof.
\par\medskip

\noindent Georg Cantor (1891). \"Uber eine elementare Frage der Mannigfaltigkeitslehre. Jahresbericht der Deutschen Mathematiker-Vereinigung 1, pp. 75--78. English: On an elementary question in the theory of manifolds, translated by Anthony Phillips, \href{https://math.stonybrook.edu/~tony/archive/336f06/cantor/cantor.html}{translation}. The diagonal argument, with the symbols m and w and the comparison with the 1874 proof on p.~75.
\par\medskip

\noindent Solomon Feferman (2006). Are There Absolutely Unsolvable Problems? G\"odel's Dichotomy. Philosophia Mathematica (III) 14(2), pp. 134--152. \href{https://doi.org/10.1093/philmat/nkj003}{DOI: 10.1093/philmat/nkj003}. The Gibbs lecture, G\"odel's private conviction, and why he stated only the either-or.
\par\medskip

\noindent Solomon Feferman (2006). The nature and significance of G\"odel's incompleteness theorems. Lecture, Institute for Advanced Study, G\"odel Centenary Program, 17 November 2006. \href{https://math.stanford.edu/~feferman/papers/Godel-IAS.pdf}{Author's text}. The announcement at K\"onigsberg met with general incomprehension except from von Neumann, and von Neumann's letter of 20 November 1930.
\par\medskip

\noindent Kurt G\"odel (1951). Some basic theorems on the foundations of mathematics and their implications. Gibbs Lecture, American Mathematical Society, Brown University, 26 December 1951. In Kurt G\"odel, Collected Works, Vol. III: Unpublished Essays and Lectures, S. Feferman et al. (eds.), Oxford University Press, 1995, pp. 304--323.
\par\medskip

\noindent Douglas R. Hofstadter (1979). G\"odel, Escher, Bach: An Eternal Golden Braid. Basic Books, New York. Where the word ``quine'' comes from.
\par\medskip

\noindent Paul A. Karger and Roger R. Schell (1974). Multics Security Evaluation: Vulnerability Analysis. Technical Report ESD-TR-74-193, Vol. II, Electronic Systems Division, Air Force Systems Command, June 1974. \href{https://apps.dtic.mil/sti/pdfs/ADA001120.pdf}{DTIC copy}. A trap door planted in a compiler that keeps itself alive.
\par\medskip

\noindent Stephen Cole Kleene (1938). On notation for ordinal numbers. Journal of Symbolic Logic 3(4), pp. 150--155. \href{https://doi.org/10.2307/2267778}{DOI: 10.2307/2267778}. The second recursion theorem, stated and proved at the end of section 2.
\par\medskip

\noindent F. William Lawvere (1969). Diagonal arguments and cartesian closed categories. In Category Theory, Homology Theory and their Applications II, Lecture Notes in Mathematics 92, Springer, pp. 134--145. \href{https://doi.org/10.1007/BFb0080769}{DOI: 10.1007/BFb0080769}. Reprinted with author commentary in Reprints in Theory and Applications of Categories 15 (2006), pp. 1--13, \href{https://www.tac.mta.ca/tac/reprints/articles/15/tr15abs.html}{reprint}.
\par\medskip

\noindent Martin H. L\"ob (1955). Solution of a problem of Leon Henkin. Journal of Symbolic Logic 20(2), pp. 115--118. \href{https://doi.org/10.2307/2266895}{DOI: 10.2307/2266895}.
\par\medskip

\noindent J. R. Lucas (1961). Minds, Machines and G\"odel. Philosophy 36(137), pp. 112--127. \href{https://doi.org/10.1017/S0031819100057983}{DOI: 10.1017/S0031819100057983}.
\par\medskip

\noindent Yiannis N. Moschovakis (2010). Kleene's Amazing Second Recursion Theorem. Bulletin of Symbolic Logic 16(2), pp. 189--239. \href{https://doi.org/10.2178/bsl/1286889124}{DOI: 10.2178/bsl/1286889124}. Extended abstract in Computer Science Logic 2009, Lecture Notes in Computer Science 5771, pp. 24--39. ``This little gem is stated unbilled and proved (completely) in the last two lines of \S2 of the short note Kleene [1938].''
\par\medskip

\noindent Caspar Oesterheld (2019). Robust program equilibrium. Theory and Decision 86(1), pp. 143--159. \href{https://doi.org/10.1007/s11238-018-9679-3}{DOI: 10.1007/s11238-018-9679-3}. Programs that cooperate by simulating each other, with a grounding step to stop the regress.
\par\medskip

\noindent Roger Penrose (1989). The Emperor's New Mind: Concerning Computers, Minds, and the Laws of Physics. Oxford University Press.
\par\medskip

\noindent Roger Penrose (1994). Shadows of the Mind: A Search for the Missing Science of Consciousness. Oxford University Press.
\par\medskip

\noindent W. V. Quine (1962). Paradox. Scientific American 206(4), pp. 84--96. \href{https://doi.org/10.1038/scientificamerican0462-84}{DOI: 10.1038/scientificamerican0462-84}. Reprinted as ``The Ways of Paradox'' in The Ways of Paradox and Other Essays (1966). The sentence ``yields falsehood when preceded by its quotation.''
\par\medskip

\noindent Panu Raatikainen (2013, revised 2025). G\"odel's Incompleteness Theorems. Stanford Encyclopedia of Philosophy. \href{https://plato.stanford.edu/entries/goedel-incompleteness/}{Entry}. Why the G\"odel sentence doesn't literally say anything about itself, and why the theorem isn't about provability in any absolute sense.
\par\medskip

\noindent Jules Richard (1905). Les principes des math\'ematiques et le probl\`eme des ensembles. Revue g\'en\'erale des sciences pures et appliqu\'ees 16(12), pp. 541--543. English translation in Jean van Heijenoort (ed.), From Frege to G\"odel, Harvard University Press, 1967, pp. 142--144. His school posts are from J. J. O'Connor and E. F. Robertson, ``Jules Antoine Richard,'' MacTutor History of Mathematics archive, \href{https://mathshistory.st-andrews.ac.uk/Biographies/Richard_Jules/}{mathshistory.st-andrews.ac.uk}.
\par\medskip

\noindent J. Barkley Rosser (1936). Extensions of some theorems of G\"odel and Church. Journal of Symbolic Logic 1(3), pp. 87--91. \href{https://doi.org/10.2307/2269028}{DOI: 10.2307/2269028}. The improvement that needs only that the system doesn't contradict itself. Rosser as a student of Alonzo Church: Wikipedia, \href{https://en.wikipedia.org/wiki/J._Barkley_Rosser}{J. Barkley Rosser}.
\par\medskip

\noindent Bertrand Russell (1903). The Principles of Mathematics. Cambridge University Press. Chapter X, The Contradiction, \S100: ``I was led to it in the endeavour to reconcile Cantor's proof that there can be no greatest cardinal number with the very plausible supposition that the class of all terms \ldots\ has necessarily the greatest possible number of members.'' \href{https://fair-use.org/bertrand-russell/the-principles-of-mathematics/s100}{Text of \S100}.
\par\medskip

\noindent Bertrand Russell (1906). Les paradoxes de la logique. Revue de m\'etaphysique et de morale 14, pp. 627--650. The first printing of Berry's paradox, credited to G. G. Berry; Berry as a junior librarian at the Bodleian: Wikipedia, \href{https://en.wikipedia.org/wiki/Berry_paradox}{Berry paradox}.
\par\medskip

\noindent Ken Thompson (1984). Reflections on Trusting Trust. Communications of the ACM 27(8), pp. 761--763. \href{https://doi.org/10.1145/358198.358210}{DOI: 10.1145/358198.358210}. Delivered as the 1983 Turing Award lecture.
\par\medskip

\noindent David A. Wheeler (2005). Countering Trusting Trust through Diverse Double-Compiling. Proceedings of the 21st Annual Computer Security Applications Conference (ACSAC 2005), Tucson, pp. 33--48. \href{https://doi.org/10.1109/CSAC.2005.17}{DOI: 10.1109/CSAC.2005.17}.
\par\medskip

\noindent David A. Wheeler (2009). Fully Countering Trusting Trust through Diverse Double-Compiling. PhD dissertation, George Mason University. \href{https://dwheeler.com/trusting-trust/}{Author's page}.
\par\medskip

\noindent Alfred North Whitehead and Bertrand Russell (1910). Principia Mathematica, Volume I. Cambridge University Press. Introduction, Chapter II, section VIII, pp. 63--64: ``the least integer not nameable in fewer than nineteen syllables.'' \href{https://www.gutenberg.org/ebooks/78050}{Project Gutenberg text}.
\par\medskip

\noindent Noson S. Yanofsky (2003). A Universal Approach to Self-Referential Paradoxes, Incompleteness and Fixed Points. Bulletin of Symbolic Logic 9(3), pp. 362--386. \href{https://doi.org/10.2178/bsl/1058448677}{DOI: 10.2178/bsl/1058448677}; \href{https://arxiv.org/abs/math/0305282}{arXiv: math/0305282}.
\par\medskip

\noindent Eliezer Yudkowsky and Marcello Herreshoff (2013). Tiling agents for self-modifying AI, and the L\"obian obstacle. Draft report, Machine Intelligence Research Institute. \href{https://intelligence.org/files/TilingAgentsDraft.pdf}{Report}.
\par\medskip

\subsection*{Chapter 18. Five real systems}

\noindent Will Arthur and David Challener, with Kenneth Goldman (2015). A Practical Guide to TPM 2.0: Using the Trusted Platform Module in the New Age of Security. Apress. Chapter 11, NV Indexes, pp. 137--150 (bit-field indexes and counters); Chapter 12, Platform Configuration Registers, pp. 151--161 (read, extend and reset). \href{https://doi.org/10.1007/978-1-4302-6584-9}{DOI: 10.1007/978-1-4302-6584-9}.
\par\medskip

\noindent Vitalik Buterin (2016). EIP-150: Gas cost changes for IO-heavy operations. Ethereum Improvement Proposals, created 24 September 2016. \href{https://eips.ethereum.org/EIPS/eip-150}{eips.ethereum.org/EIPS/eip-150}. The repair switched on in the Tangerine Whistle fork, raising the prices of the opcodes that read the state from disk.
\par\medskip

\noindent Vitalik Buterin (2016). Transaction spam attack: next steps. Ethereum Foundation Blog, 22 September 2016. \href{https://blog.ethereum.org/2016/09/22/transaction-spam-attack-next-steps/}{Post}. The 2016 attack, which made every node read state from disk with a cheap opcode.
\par\medskip

\noindent Vitalik Buterin and Virgil Griffith (2017). Casper the Friendly Finality Gadget. \href{https://arxiv.org/abs/1710.09437v1}{arXiv: 1710.09437v1}. Justified and final checkpoints, the two slashing rules, and accountable safety, for a finality system laid over a proof-of-work chain.
\par\medskip

\noindent Vitalik Buterin and Martin Swende (2020). EIP-2929: Gas cost increases for state access opcodes. Ethereum Improvement Proposals. \href{https://eips.ethereum.org/EIPS/eip-2929}{eips.ethereum.org/EIPS/eip-2929}. The cold-slot charge of 2,100, justified by the 2016 attacks and switched on in the Berlin upgrade of April 2021.
\par\medskip

\noindent Vitalik Buterin, Diego Hernandez, Thor Kamphefner, Khiem Pham, Zhi Qiao, Danny Ryan, Juhyeok Sin, Ying Wang and Yan X. Zhang (2020). Combining GHOST and Casper. \href{https://arxiv.org/abs/2003.03052}{arXiv: 2003.03052}. Gasper, which joins Casper's finality to the LMD GHOST rule for picking branches.
\par\medskip

\noindent Byung-Gon Chun, Petros Maniatis, Scott Shenker and John Kubiatowicz (2007). Attested append-only memory: making adversaries stick to their word. Proceedings of the 21st ACM Symposium on Operating Systems Principles (SOSP '07), pp. 189--204. \href{https://doi.org/10.1145/1294261.1294280}{DOI: 10.1145/1294261.1294280}.
\par\medskip

\noindent Allen Clement, Flavio Junqueira, Aniket Kate and Rodrigo Rodrigues (2012). On the (limited) power of non-equivocation. Proceedings of the 2012 ACM Symposium on Principles of Distributed Computing (PODC '12), pp. 301--308. \href{https://doi.org/10.1145/2332432.2332490}{DOI: 10.1145/2332432.2332490}. Stopping a copy from saying two different things doesn't lower the 3f + 1 bound on its own, and with signatures a third party can check it turns a crash-tolerant protocol into one that tolerates liars with no extra copies.
\par\medskip

\noindent Victor Costan and Srinivas Devadas (2016). Intel SGX Explained. Cryptology ePrint Archive, Paper 2016/086. \href{https://eprint.iacr.org/2016/086}{IACR ePrint 2016/086}. How SGX enclaves are measured and how they produce reports of their own, apart from the chip's registers.
\par\medskip

\noindent Scott A. Crosby and Dan S. Wallach (2009). Efficient data structures for tamper-evident logging. Proceedings of the 18th USENIX Security Symposium, Montreal. \href{https://www.usenix.org/conference/usenixsecurity09/technical-sessions/presentation/efficient-data-structures-tamper-evident}{Conference page}. Logs built on hash trees at Rice University, with short proofs that an entry is in and that a later version extends an earlier one.
\par\medskip

\noindent Joe DeBlasio, for the Chrome CT team (2025). Discussion: removal of one-6962-log requirement. Certificate Transparency Policy group, 17 November 2025. \href{https://groups.google.com/a/chromium.org/g/ct-policy/c/amrnTiDFdwM/m/mCyPDc4TAwAJ}{Post}. The proposal to stop requiring every certificate to carry a timestamp from at least one RFC 6962 log.
\par\medskip

\noindent J\'er\'emie Decouchant, David Kozhaya, Vincent Rahli and Jiangshan Yu (2022). Damysus: streamlined BFT consensus leveraging trusted components. Proceedings of the Seventeenth European Conference on Computer Systems (EuroSys '22), pp. 1--16. \href{https://doi.org/10.1145/3492321.3519568}{DOI: 10.1145/3492321.3519568}. The small trusted service called the Checker.
\par\medskip

\noindent Devcon (n.d.). Past events. \href{https://devcon.org/en/past-events/}{Page}. Devcon 0 in Berlin, Devcon 1 in London and Devcon 2 in Shanghai, where denial-of-service attacks hit the network hours before the conference began.
\par\medskip

\noindent Ethos (2016). Ethereum's Shanghai attacks. ethos.dev, 30 November 2016. \href{https://ethos.dev/shanghai-attacks}{Article}. The attacks that began during the second developer conference in Shanghai, and the October 2016 fork that repriced the cheap opcodes.
\par\medskip

\noindent Ethereum Foundation (n.d.). Ethereum execution layer specifications (execution-specs). GitHub repository. \href{https://github.com/ethereum/execution-specs}{Repository}. The specification the copied gas and refund arithmetic for one storage slot follows.
\par\medskip

\noindent Ethereum Foundation (n.d.). Timeline of all Ethereum forks. ethereum.org. \href{https://ethereum.org/en/ethereum-forks/}{Page}. Tangerine Whistle on 18 October 2016, Berlin on 15 April 2021, London on 5 August 2021, and the Merge to proof of stake on 15 September 2022.
\par\medskip

\noindent Yoichi Hirai (2017). Isabelle/HOL proofs of accountable safety for one-message Casper, in the repository pos. \href{https://github.com/palmskog/pos}{Repository}.
\par\medskip

\noindent Ben Laurie, Adam Langley and Emilia Kasper (2013). Certificate Transparency. RFC 6962, Internet Engineering Task Force, June 2013. \href{https://doi.org/10.17487/RFC6962}{DOI: 10.17487/RFC6962}. The version the running logs and browsers still follow.
\par\medskip

\noindent Ben Laurie, Eran Messeri and Rob Stradling (2021). Certificate Transparency Version 2.0. RFC 9162, Internet Engineering Task Force, December 2021. \href{https://doi.org/10.17487/RFC9162}{DOI: 10.17487/RFC9162}. The version copied here, which says logs let misissued certificates be detected afterwards without preventing them, and lists a log showing different trees to different parties among the ways a log can misbehave.
\par\medskip

\noindent Dave Levin, John R. Douceur, Jacob R. Lorch and Thomas Moscibroda (2009). TrInc: small trusted hardware for large distributed systems. Proceedings of the 6th USENIX Symposium on Networked Systems Design and Implementation (NSDI '09), pp. 1--14. \href{https://www.usenix.org/legacy/event/nsdi09/tech/full_papers/levin/levin.pdf}{Paper}. A trusted component that is only a counter that never goes down and a key.
\par\medskip

\noindent Sinisa Matetic, Mansoor Ahmed, Kari Kostiainen, Aritra Dhar, David Sommer, Arthur Gervais, Ari Juels and Srdjan Capkun (2017). ROTE: rollback protection for trusted execution. Proceedings of the 26th USENIX Security Symposium, pp. 1289--1306. \href{https://www.usenix.org/conference/usenixsecurity17/technical-sessions/presentation/matetic}{USENIX page}. Several machines vouch for each other's counters, so the only way left to break integrity is to reset all of them to their initial state.
\par\medskip

\noindent Mozilla (n.d.). Certificate Transparency. MDN Web Docs. \href{https://developer.mozilla.org/en-US/docs/Web/Security/Defenses/Certificate_Transparency}{Page}. Browser implementations are based on RFC 6962, though RFC 9162 is the current specification.
\par\medskip

\noindent Devon O'Brien (2018). Certificate Transparency Enforcement in Google Chrome. Certificate Transparency Policy group, 6 February 2018. \href{https://groups.google.com/a/chromium.org/g/ct-policy/c/wHILiYf31DE/m/iMFmpMEkAQAJ}{Post}. Certificates issued after 30 April 2018 must comply with Chrome's CT policy or are not trusted, with a full-page warning for the visitor.
\par\medskip

\noindent Karl Palmskog, Milos Gligoric, Lucas Pe\~na, Brandon Moore and Grigore Ro\c{s}u (2018). Verification of Casper in the Coq Proof Assistant. Technical report, Runtime Verification, Inc., The University of Texas at Austin and University of Illinois at Urbana-Champaign. \href{https://github.com/runtimeverification/casper-proofs/blob/master/report/report.pdf}{Report}.
\par\medskip

\noindent Bryan Parno, Jacob R. Lorch, John R. Douceur, James Mickens and Jonathan M. McCune (2011). Memoir: practical state continuity for protected modules. Proceedings of the 2011 IEEE Symposium on Security and Privacy, pp. 379--394. \href{https://doi.org/10.1109/SP.2011.38}{DOI: 10.1109/SP.2011.38}. The counter that need only tick once every five seconds, the storage rated for about 100,000 writes, and the password checker rolled back to undo a password change.
\par\medskip

\noindent Rob Smart (2021). Implementing the 10 Security Goals: Anti-rollback Explained. PSA Certified, 27 April 2021. \href{https://www.psacertified.org/blog/anti-rollback-explained/}{Article}. A counter kept in one-time-programmable fuse bits, with enough bits for the updates expected over the root of trust's supported lifetime.
\par\medskip

\noindent Raoul Strackx and Frank Piessens (2016). Ariadne: a minimal approach to state continuity. Proceedings of the 25th USENIX Security Symposium, pp. 875--892. \href{https://www.usenix.org/conference/usenixsecurity16/technical-sessions/presentation/strackx}{USENIX page}. A single bit flip of lasting memory per state update, which the authors call the theoretical lower limit.
\par\medskip

\noindent Runtime Verification, Inc. and The University of Texas at Austin (2018). casper-proofs: models and proofs of Casper FFG accountable safety and plausible liveness. \href{https://github.com/runtimeverification/casper-proofs}{Repository}. The accountable-safety argument copied here, built on Yoichi Hirai's proof with the two-thirds overlap taken as a condition.
\par\medskip

\noindent Giuliana Santos Veronese, Miguel Correia, Alysson Neves Bessani, Lau Cheuk Lung and Paulo Verissimo (2013). Efficient Byzantine fault-tolerance. IEEE Transactions on Computers 62(1), pp. 16--30. \href{https://doi.org/10.1109/TC.2011.221}{DOI: 10.1109/TC.2011.221}. MinBFT, which uses a trusted counter to agree with 2f + 1 copies.
\par\medskip

\subsection*{Chapter 19. Which box is worth keeping}

\noindent J\'anos Acz\'el, Bruno Forte and Che Tat Ng (1974). Why the Shannon and Hartley entropies are ``natural.'' Advances in Applied Probability 6(1), pp. 131--146. \href{https://doi.org/10.2307/1426210}{DOI: 10.2307/1426210}. The entropies with a short list of natural properties are exactly the mixtures of Shannon's and Hartley's.
\par\medskip

\noindent John C. Baez, Tobias Fritz and Tom Leinster (2011). A characterization of entropy in terms of information loss. Entropy 13(11), pp. 1945--1957. \href{https://doi.org/10.3390/e13111945}{DOI: 10.3390/e13111945}; \href{https://arxiv.org/abs/1106.1791}{arXiv: 1106.1791}. Shannon's entropy, up to a constant, as the only measure of information loss that adds along a chain of maps, mixes the way probabilities mix, and changes smoothly.
\par\medskip

\noindent Philip Ball (2019). Quantum Darwinism, an Idea to Explain Objective Reality, Passes First Tests. Quanta Magazine, 22 July 2019. \href{https://www.quantamagazine.org/quantum-darwinism-an-idea-to-explain-objective-reality-passes-first-tests-20190722/}{Article}. The three experiments in Italy, China and Germany, with Zurek's and others' cautions about them.
\par\medskip

\noindent Fernando G. S. L. Brand\~ao, Marco Piani and Pawe\l{} Horodecki (2015). Generic emergence of classical features in quantum Darwinism. Nature Communications 6, 7908. \href{https://doi.org/10.1038/ncomms8908}{DOI: 10.1038/ncomms8908}.
\par\medskip

\noindent Vitalik Buterin, Eric Conner, Rick Dudley, Matthew Slipper, Ian Norden and Abdelhamid Bakhta (2019). EIP-1559: Fee market change for ETH 1.0 chain. Ethereum Improvement Proposals, created 13 April 2019. \href{https://eips.ethereum.org/EIPS/eip-1559}{eips.ethereum.org/EIPS/eip-1559}. The base fee aimed at a half-full block, moving by at most an eighth per block.
\par\medskip

\noindent Ming-Cheng Chen, Han-Sen Zhong, Yuan Li, Dian Wu, Xi-Lin Wang, Li Li, Nai-Le Liu, Chao-Yang Lu and Jian-Wei Pan (2019). Emergence of classical objectivity of quantum Darwinism in a photonic quantum simulator. Science Bulletin 64(9), pp. 580--585. \href{https://doi.org/10.1016/j.scib.2019.03.032}{DOI: 10.1016/j.scib.2019.03.032}. The experiment in China, done with light.
\par\medskip

\noindent Mario A. Ciampini, Giorgia Pinna, Paolo Mataloni and Mauro Paternostro (2018). Experimental signature of quantum Darwinism in photonic cluster states. Physical Review A 98, 020101. \href{https://doi.org/10.1103/PhysRevA.98.020101}{DOI: 10.1103/PhysRevA.98.020101}. The experiment in Italy, done with light.
\par\medskip

\noindent Philippe Allard Gu\'erin, Veronika Baumann, Flavio Del Santo and \v{C}aslav Brukner (2021). A no-go theorem for the persistent reality of Wigner's friend's perception. Communications Physics 4, 93. \href{https://doi.org/10.1038/s42005-021-00589-1}{DOI: 10.1038/s42005-021-00589-1}.
\par\medskip

\noindent J. J. Halliwell (1999). Somewhere in the universe: Where is the information stored when histories decohere? Physical Review D 60, 105031. \href{https://doi.org/10.1103/PhysRevD.60.105031}{DOI: 10.1103/PhysRevD.60.105031}. The link between records and decoherent histories carried from pure starting states to mixed ones, where records can still exist but may be imperfect.
\par\medskip

\noindent Omer Tamuz, Caltech, \href{https://www.hss.caltech.edu/people/omer-tamuz}{faculty page}; Philipp Strack, Yale, \href{https://economics.yale.edu/people/philipp-strack}{faculty page}; Luciano Pomatto, Caltech, \href{https://hss.caltech.edu/documents/3303/HSS_New_Faces_2016-2017_FINAL.pdf}{HSS new faculty 2016--17}. The economists behind the 2023 cost-of-information paper.
\par\medskip

\noindent Luciano Pomatto, Philipp Strack and Omer Tamuz (2023). The cost of information: the case of constant marginal costs. American Economic Review 113(5), pp. 1360--1393. \href{https://doi.org/10.1257/aer.20190185}{DOI: 10.1257/aer.20190185}. A uniqueness result for the cost of getting information when running an experiment twice costs twice as much.
\par\medskip

\noindent Xiao-Liang Qi and Daniel Ranard (2021). Emergent classicality in general multipartite states and channels. Quantum 5, article 555. \href{https://doi.org/10.22331/q-2021-09-28-555}{DOI: 10.22331/q-2021-09-28-555}.
\par\medskip

\noindent C. Jess Riedel, Wojciech H. Zurek and Michael Zwolak (2016). Objective past of a quantum universe: Redundant records of consistent histories. Physical Review A 93, 032126. \href{https://doi.org/10.1103/PhysRevA.93.032126}{DOI: 10.1103/PhysRevA.93.032126}; \href{https://arxiv.org/abs/1312.0331}{arXiv: 1312.0331}.
\par\medskip

\noindent Tom Leinster. Wikipedia. \href{https://en.wikipedia.org/wiki/Tom_Leinster}{Article}. A British mathematician known for his work on category theory, at Glasgow and then Edinburgh. John C. Baez, Wikipedia, \href{https://en.wikipedia.org/wiki/John_C._Baez}{Article}: mathematical physicist at the University of California, Riverside, working on higher categories and applied category theory. Tobias Fritz's category-theoretic approach to probability: \href{https://arxiv.org/a/fritz_t_1}{arXiv listing}.
\par\medskip

\noindent T. K. Unden, D. Louzon, M. Zwolak, W. H. Zurek and F. Jelezko (2019). Revealing the Emergence of Classicality Using Nitrogen-Vacancy Centers. Physical Review Letters 123, 140402. \href{https://doi.org/10.1103/PhysRevLett.123.140402}{DOI: 10.1103/PhysRevLett.123.140402}. The experiment at Ulm on a single defect in diamond and the carbon atoms around it.
\par\medskip

\noindent University of Waterloo, Department of Pure Mathematics (2020). We remember Distinguished Professor Emeritus János Aczél. \href{https://uwaterloo.ca/pure-mathematics/news/we-remember-distinguished-professor-emeritus-janos-aczel}{Notice}. Aczél at Waterloo from 1965, founder of Aequationes Mathematicae. Bruno Forte at Waterloo: \href{https://mathgenealogy.com/id.php?id=30611}{Mathematics Genealogy Project}. Che Tat Ng's doctorate, Waterloo 1971, under Aczél: \href{https://mathgenealogy.org/id.php?id=102040}{Mathematics Genealogy Project}.
\par\medskip

\noindent Wojciech H. Zurek (1981). Pointer basis of quantum apparatus: Into what mixture does the wave packet collapse? Physical Review D 24(6), pp. 1516--1525. \href{https://doi.org/10.1103/PhysRevD.24.1516}{DOI: 10.1103/PhysRevD.24.1516}.
\par\medskip

\noindent Wojciech H. Zurek (2009). Quantum Darwinism. Nature Physics 5(3), pp. 181--188. \href{https://doi.org/10.1038/nphys1202}{DOI: 10.1038/nphys1202}; \href{https://arxiv.org/abs/0903.5082v1}{arXiv: 0903.5082v1}.
\par\medskip

\subsection*{Chapter 20. Heat}

\noindent Erik Aurell, Carlos Mej\'ia-Monasterio and Paolo Muratore-Ginanneschi (2011). Optimal protocols and optimal transport in stochastic thermodynamics. Physical Review Letters 106, 250601. \href{https://doi.org/10.1103/PhysRevLett.106.250601}{DOI: 10.1103/PhysRevLett.106.250601}; \href{https://arxiv.org/abs/1012.2037}{arXiv: 1012.2037}.
\par\medskip

\noindent Antoine B\'erut, Artyom Petrosyan and Sergio Ciliberto (2013). Detailed Jarzynski equality applied to a logically irreversible procedure. EPL (Europhysics Letters) 103(6), 60002. \href{https://doi.org/10.1209/0295-5075/103/60002}{DOI: 10.1209/0295-5075/103/60002}; \href{https://arxiv.org/abs/1302.4417}{arXiv: 1302.4417}. The bound of kT times ln 2 plus ln of the success rate on the work averaged over the successful erasures.
\par\medskip

\noindent Antoine B\'erut, Artyom Petrosyan and Sergio Ciliberto (2015). Information and thermodynamics: experimental verification of Landauer's erasure principle. Journal of Statistical Mechanics: Theory and Experiment, P06015. \href{https://doi.org/10.1088/1742-5468/2015/06/P06015}{DOI: 10.1088/1742-5468/2015/06/P06015}; \href{https://arxiv.org/abs/1503.06537}{arXiv: 1503.06537}. The long version, with the protocol in full, the fit $\langle Q\rangle = kT\ln 2 + B/\tau$ with $B = 8.15\,kT\cdot\mathrm{s}$, and the hopping times of about 3,000 seconds above an $8\,kT$ barrier and about ten seconds at $2.2\,kT$.
\par\medskip

\noindent John Earman and John D. Norton (1998, 1999). Exorcist XIV: The Wrath of Maxwell's Demon. Part I: From Maxwell to Szilard, Studies in History and Philosophy of Modern Physics 29(4), pp. 435--471, \href{https://doi.org/10.1016/S1355-2198(98)00023-9}{DOI: 10.1016/S1355-2198(98)00023-9}; Part II: From Szilard to Landauer and Beyond, 30(1), pp. 1--40, \href{https://doi.org/10.1016/S1355-2198(98)00026-4}{DOI: 10.1016/S1355-2198(98)00026-4}. Author copies: \href{https://sites.pitt.edu/~jdnorton/papers/ExorcistXIV/Exorcist1.pdf}{Part I}; \href{https://sites.pitt.edu/~jdnorton/papers/ExorcistXIV/Exorcist2.pdf}{Part II}. The dilemma with its sound and profound horns is set out at the start of Part II.
\par\medskip

\noindent Rocco Gaudenzi, Enrique Burzur\'i, Satoru Maegawa, Herre S. J. van der Zant and Fernando Luis (2018). Quantum Landauer erasure with a molecular nanomagnet. Nature Physics 14, pp. 565--568. \href{https://doi.org/10.1038/s41567-018-0070-7}{DOI: 10.1038/s41567-018-0070-7}. Landauer's floor on a register of molecular spins at one kelvin.
\par\medskip

\noindent Mom\v{c}ilo Gavrilov and John Bechhoefer (2016). Erasure without Work in an Asymmetric Double-Well Potential. Physical Review Letters 117, 200601. \href{https://doi.org/10.1103/PhysRevLett.117.200601}{DOI: 10.1103/PhysRevLett.117.200601}. Average erasure work below $kT\ln 2$ when the two wells are not alike.
\par\medskip

\noindent Anson Ho, Ege Erdil and Tamay Besiroglu (2023). Limits to the Energy Efficiency of CMOS Microprocessors. \href{https://arxiv.org/abs/2312.08595}{arXiv: 2312.08595}. Table I: switching a CMOS gate, 100 aJ to 100 fJ.
\par\medskip

\noindent Jeongmin Hong, Brian Lambson, Scott Dhuey and Jeffrey Bokor (2016). Experimental test of Landauer's principle in single-bit operations on nanomagnetic memory bits. Science Advances 2(3), e1501492. \href{https://doi.org/10.1126/sciadv.1501492}{DOI: 10.1126/sciadv.1501492}. The floor on a real memory technology, $1.45 \pm 0.35\,kT$, within two standard deviations.
\par\medskip

\noindent Christopher Jarzynski (1997). Nonequilibrium Equality for Free Energy Differences. Physical Review Letters 78(14), pp. 2690--2693. \href{https://doi.org/10.1103/PhysRevLett.78.2690}{DOI: 10.1103/PhysRevLett.78.2690}.
\par\medskip

\noindent Ryoichi Kawai, Juan M. R. Parrondo and Christian Van den Broeck (2007). Dissipation: the phase-space perspective. Physical Review Letters 98, 080602. \href{https://doi.org/10.1103/PhysRevLett.98.080602}{DOI: 10.1103/PhysRevLett.98.080602}. The detailed version of Jarzynski's equality for a partition of the states.
\par\medskip

\noindent Artemy Kolchinsky and David H. Wolpert (2017). Dependence of dissipation on the initial distribution over states. Journal of Statistical Mechanics: Theory and Experiment 2017, 083202. \href{https://doi.org/10.1088/1742-5468/aa7ee1}{DOI: 10.1088/1742-5468/aa7ee1}. The extra cost of running a fixed process on a spread it wasn't built for, which they call dissipation due to incorrect priors.
\par\medskip

\noindent James Ladyman, Stuart Presnell, Anthony J. Short and Berry Groisman (2007). The connection between logical and thermodynamic irreversibility. Studies in History and Philosophy of Modern Physics 38(1), pp. 58--79. \href{https://doi.org/10.1016/j.shpsb.2006.03.007}{DOI: 10.1016/j.shpsb.2006.03.007}. From Kelvin's statement of the second law, no L-machine carries out a logically irreversible operation in a thermodynamically reversible way.
\par\medskip

\noindent Gonzalo Manzano, G\"ulce Karde\c{s}, \'Edgar Rold\'an and David H. Wolpert (2024). Thermodynamics of Computations with Absolute Irreversibility, Unidirectional Transitions, and Stochastic Computation Times. Physical Review X 14, 021026. \href{https://doi.org/10.1103/PhysRevX.14.021026}{DOI: 10.1103/PhysRevX.14.021026}.
\par\medskip

\noindent O. J. E. Maroney (2009). Generalizing Landauer's Principle. Physical Review E 79, 031105. \href{https://doi.org/10.1103/PhysRevE.79.031105}{DOI: 10.1103/PhysRevE.79.031105}; \href{https://arxiv.org/abs/quant-ph/0702094v3}{arXiv: quant-ph/0702094v3}. Written at the Perimeter Institute; the general form of the principle, and the floor for a memory whose two states differ.
\par\medskip

\noindent Owen Maroney (2009). Information Processing and Thermodynamic Entropy. Stanford Encyclopedia of Philosophy, first published 15 September 2009. \href{https://plato.stanford.edu/entries/information-entropy/}{Entry}. It ends without a verdict on Landauer's principle.
\par\medskip

\noindent Y\^uto Murashita, Ken Funo and Masahito Ueda (2014). Nonequilibrium equalities in absolutely irreversible processes. Physical Review E 90, 042110. \href{https://doi.org/10.1103/PhysRevE.90.042110}{DOI: 10.1103/PhysRevE.90.042110}; \href{https://arxiv.org/abs/1401.4494}{arXiv: 1401.4494}. Absolute irreversibility, and the free expansion of a gas from one part of a box split in the ratio l to 1 minus l.
\par\medskip

\noindent John D. Norton (2005). Eaters of the Lotus: Landauer's Principle and the Return of Maxwell's Demon. Studies in History and Philosophy of Modern Physics 36(2), pp. 375--411. \href{https://doi.org/10.1016/j.shpsb.2004.12.002}{DOI: 10.1016/j.shpsb.2004.12.002}.
\par\medskip

\noindent John D. Norton (2018). Maxwell's demon does not compute. In M. E. Cuffaro and S. C. Fletcher (eds.), Physical Perspectives on Computation, Computational Perspectives on Physics, Cambridge University Press, pp. 240--256. \href{https://doi.org/10.1017/9781316759745.012}{DOI: 10.1017/9781316759745.012}; \href{https://philsci-archive.pitt.edu/14808/}{PhilSci-Archive 14808}. The demon exorcised with Liouville's theorem alone.
\par\medskip

\noindent John D. Norton (2026). Too good to be true: entropy, information, computation, and the uninformed thermodynamics of erasure. Philosophy of Science, published online 22 May 2026; manuscript dated 2025. \href{https://doi.org/10.1017/psa.2026.10221}{DOI: 10.1017/psa.2026.10221}. The careful literature ``almost invariably accepts the sound horn,'' and what works there is relabelled ordinary thermodynamics.
\par\medskip

\noindent Juan M. R. Parrondo, Jordan M. Horowitz and Takahiro Sagawa (2015). Thermodynamics of information. Nature Physics 11(2), pp. 131--139. \href{https://doi.org/10.1038/nphys3230}{DOI: 10.1038/nphys3230}.
\par\medskip

\noindent Paul M. Riechers, Alexander B. Boyd, Gregory W. Wimsatt and James P. Crutchfield (2020). Balancing error and dissipation in computing. Physical Review Research 2, 033524. \href{https://doi.org/10.1103/PhysRevResearch.2.033524}{DOI: 10.1103/PhysRevResearch.2.033524}; \href{https://arxiv.org/abs/1909.06650}{arXiv: 1909.06650}. Time-symmetric control and nonreciprocal transitions.
\par\medskip

\noindent Takahiro Sagawa (2014). Thermodynamic and logical reversibilities revisited. Journal of Statistical Mechanics: Theory and Experiment, P03025. \href{https://doi.org/10.1088/1742-5468/2014/03/P03025}{DOI: 10.1088/1742-5468/2014/03/P03025}; \href{https://arxiv.org/abs/1311.1886v2}{arXiv: 1311.1886v2}. Erasure that can in principle approach producing no entropy overall.
\par\medskip

\noindent Tim Schmiedl and Udo Seifert (2007). Optimal finite-time processes in stochastic thermodynamics. Physical Review Letters 98, 108301. \href{https://doi.org/10.1103/PhysRevLett.98.108301}{DOI: 10.1103/PhysRevLett.98.108301}; \href{https://arxiv.org/abs/cond-mat/0701554}{arXiv: cond-mat/0701554}. The least work to drag a trapped bead a set distance in a set time.
\par\medskip

\noindent J. Schnakenberg (1976). Network theory of microscopic and macroscopic behavior of master equation systems. Reviews of Modern Physics 48(4), pp. 571--585. \href{https://doi.org/10.1103/RevModPhys.48.571}{DOI: 10.1103/RevModPhys.48.571}. The entropy produced by a system that jumps between states.
\par\medskip

\noindent Udo Seifert (2012). Stochastic thermodynamics, fluctuation theorems and molecular machines. Reports on Progress in Physics 75(12), 126001. \href{https://doi.org/10.1088/0034-4885/75/12/126001}{DOI: 10.1088/0034-4885/75/12/126001}; \href{https://arxiv.org/abs/1205.4176}{arXiv: 1205.4176}. The local-detailed-balance premise linking jump ratios to heat.
\par\medskip

\noindent Claude E. Shannon, interviewed by Robert Price (1982). Oral history, 28 July 1982, IEEE History Center. \href{https://ethw.org/Oral-History:Claude_E._Shannon}{Transcript}. Excerpts printed as A conversation with Claude Shannon, IEEE Communications Magazine 22(5), pp. 123--126, May 1984, \href{https://doi.org/10.1109/MCOM.1984.1091957}{DOI: 10.1109/MCOM.1984.1091957}. Asked whether von Neumann gave him the word ``entropy,'' Shannon answered, ``No, I don't think he did.''
\par\medskip

\noindent Naoto Shiraishi, Ken Funo and Keiji Saito (2018). Speed limit for classical stochastic processes. Physical Review Letters 121, 070601. \href{https://doi.org/10.1103/PhysRevLett.121.070601}{DOI: 10.1103/PhysRevLett.121.070601}; \href{https://arxiv.org/abs/1802.06554}{arXiv: 1802.06554}.
\par\medskip

\noindent Myron Tribus and Edward C. McIrvine (1971). Energy and Information. Scientific American 225(3), pp. 179--188. \href{https://doi.org/10.1038/scientificamerican0971-179}{DOI: 10.1038/scientificamerican0971-179}. The story of von Neumann telling Shannon to call his measure entropy.
\par\medskip

\noindent Suriyanarayanan Vaikuntanathan and Christopher Jarzynski (2009). Dissipation and lag in irreversible processes. EPL (Europhysics Letters) 87(6), 60005. \href{https://doi.org/10.1209/0295-5075/87/60005}{DOI: 10.1209/0295-5075/87/60005}.
\par\medskip

\noindent David H. Wolpert and Artemy Kolchinsky (2020). Thermodynamics of computing with circuits. New Journal of Physics 22, 063047. \href{https://doi.org/10.1088/1367-2630/ab82b8}{DOI: 10.1088/1367-2630/ab82b8}. First posted in 2018 as Exact, complete expressions for the thermodynamic costs of circuits, \href{https://arxiv.org/abs/1806.04103v1}{arXiv: 1806.04103v1}, where the extra cost of a mismatched spread is called mismatch cost.
\par\medskip

\subsection*{Chapter 21. Keeping without the heat}

\noindent Jacob D. Bekenstein (1981). Universal upper bound on the entropy-to-energy ratio for bounded systems. Physical Review D 23(2), pp. 287--298. \href{https://doi.org/10.1103/PhysRevD.23.287}{DOI: 10.1103/PhysRevD.23.287}. The bound ties entropy to energy and size under its own physical assumptions; it isn't a measured count of a device's states.
\par\medskip

\noindent John S. Bell (1990). Against ``measurement''. Physics World 3(8), pp. 33--41. \href{https://doi.org/10.1088/2058-7058/3/8/26}{DOI: 10.1088/2058-7058/3/8/26}. The list of words with no place in a physics that wants to be precise, with measurement the worst of all.
\par\medskip

\noindent Charles H. Bennett (1989). Time/Space Trade-Offs for Reversible Computation. SIAM Journal on Computing 18(4), pp. 766--776. \href{https://doi.org/10.1137/0218053}{DOI: 10.1137/0218053}. A reversible machine saves room by undoing and redoing its own intermediate work, paying for it in time.
\par\medskip

\noindent Niels Bohr (1958). Atomic Physics and Human Knowledge. John Wiley and Sons, New York. The essay Unity of Knowledge, an address given in October 1954 for the bicentennial of Columbia University, first printed in The Unity of Knowledge (Doubleday, 1955). The line about permanent marks on a photographic plate is on p.~73, where the sentence ends with electrons entering the emulsion, not with water drops in a cloud chamber.
\par\medskip

\noindent Sandro Donadi, Kristian Piscicchia, Catalina Curceanu, Lajos Di\'osi, Matthias Laubenstein and Angelo Bassi (2021). Underground test of gravity-related wave function collapse. Nature Physics 17, pp. 74--78. \href{https://doi.org/10.1038/s41567-020-1008-4}{DOI: 10.1038/s41567-020-1008-4}. The radiation a gravity-driven collapse would give off, listened for at Gran Sasso and not found at the level the version with no adjustable settings predicts.
\par\medskip

\noindent Mauro Dorato and Jan Faye (2026). Niels Bohr as a Physicist of Principle. Preprint. \href{https://arxiv.org/abs/2609.04297}{arXiv: 2609.04297}. Where I first found Bohr's line; their quotation runs on into a clause about a water drop forming around an ion in a cloud chamber, which isn't on Bohr's p.~73.
\par\medskip

\noindent Bernhard Rauer, Sebastian Erne, Thomas Schweigler, Federica Cataldini, Mohammadamin Tajik and J\"org Schmiedmayer (2018). Recurrences in an isolated quantum many-body system. Science 360(6386), pp. 307--310. \href{https://doi.org/10.1126/science.aan7938}{DOI: 10.1126/science.aan7938}; \href{https://arxiv.org/abs/1705.08231}{arXiv: 1705.08231}. The coherence of an isolated gas of thousands of atoms came partly back around.
\par\medskip

\noindent David Reeb and Michael M. Wolf (2014). An improved Landauer principle with finite-size corrections. New Journal of Physics 16, 103011. \href{https://doi.org/10.1088/1367-2630/16/10/103011}{DOI: 10.1088/1367-2630/16/10/103011}; \href{https://arxiv.org/abs/1306.4352}{arXiv: 1306.4352}. Their heat bound assumes a reservoir that starts thermal and uncorrelated with the memory, and a combined motion that is unitary.
\par\medskip

\noindent Andrea Vinante, Rima Mezzena, Paolo Falferi, Matteo Carlesso and Angelo Bassi (2017). Improved noninterferometric test of collapse models using ultracold cantilevers. Physical Review Letters 119, 110401. \href{https://doi.org/10.1103/PhysRevLett.119.110401}{DOI: 10.1103/PhysRevLett.119.110401}. A cantilever at millikelvin with a force noise of unknown origin, compatible with the heating one collapse model predicts.
\par\medskip

\noindent John Archibald Wheeler (1989). Information, Physics, Quantum: The Search for Links. Proceedings of the 3rd International Symposium on Foundations of Quantum Mechanics, Tokyo, 1989, pp. 354--368. \href{https://jawarchive.files.wordpress.com/2012/03/informationquantumphysics.pdf}{Paper}. Also printed in Wojciech H. Zurek (ed.), Complexity, Entropy and the Physics of Information, Addison-Wesley, 1990. It from bit.
\par\medskip

\subsection*{Chapter 22. What's shown, what's chosen, what's argued}

\noindent Alok Aggarwal and Jeffrey Scott Vitter (1988). The input/output complexity of sorting and related problems. Communications of the ACM 31(9), pp. 1116--1127. \href{https://doi.org/10.1145/48529.48535}{DOI: 10.1145/48529.48535}. A model that counts only the blocks moved between a slow memory and a fast one.
\par\medskip

\noindent Pablo Arrighi and Gilles Dowek (2012). The physical Church-Turing thesis and the principles of quantum theory. International Journal of Foundations of Computer Science 23(5), pp. 1131--1145. \href{https://doi.org/10.1142/S0129054112500153}{DOI: 10.1142/S0129054112500153}; \href{https://arxiv.org/abs/1102.1612}{arXiv: 1102.1612}. Gandy's argument laid out from homogeneity of space and time, bounded density and velocity of information, and a quiet background outside a finite region.
\par\medskip

\noindent Alexander Bastounis, Fabian Circelli and Anders C. Hansen (2026). Navier-Stokes lost in translation: Why Lean verification of AI autoformalisation does not guarantee correct natural language proofs. arXiv:2610.08144. \href{https://arxiv.org/abs/2610.08144}{arXiv}. Argues that the version of OpenAI's proof written for the proof checker doesn't match the proof written in words.
\par\medskip

\noindent Tristan Buckmaster (2026). Statement on the forced blowup results with Levent Alp\"oge and on events with OpenAI. Author's page, Courant Institute, New York University. \href{https://cims.nyu.edu/~tristanb/statement.pdf}{PDF}. Forced blowup for incompressible porous media, Boussinesq and three-dimensional Euler, and his account of the dispute over credit.
\par\medskip

\noindent Stephen A. Cook and Robert A. Reckhow (1973). Time bounded random access machines. Journal of Computer and System Sciences 7(4), pp. 354--375. \href{https://doi.org/10.1016/S0022-0000(73)80029-7}{DOI: 10.1016/S0022-0000(73)80029-7}.
\par\medskip

\noindent Erik D. Demaine, Jayson Lynch, Geronimo J. Mirano and Nirvan Tyagi (2016). Energy-Efficient Algorithms. Proceedings of the 2016 ACM Conference on Innovations in Theoretical Computer Science (ITCS), pp. 321--332. \href{https://doi.org/10.1145/2840728.2840756}{DOI: 10.1145/2840728.2840756}. Energy counted beside time and space, charging for the bits an instruction destroys.
\par\medskip

\noindent Charles L. Fefferman (n.d.). Existence and Smoothness of the Navier--Stokes Equation. Official problem description, Clay Mathematics Institute. \href{https://www.claymath.org/wp-content/uploads/2022/06/navierstokes.pdf}{PDF}. The four statements a solution may prove; (C) and (D) are breakdown with a smooth force, on all of space and on the periodic box.
\par\medskip

\noindent Robin Gandy (1980). Church's Thesis and Principles for Mechanisms. In The Kleene Symposium, Studies in Logic and the Foundations of Mathematics, North-Holland, Amsterdam, pp. 123--148. \href{https://doi.org/10.1016/S0049-237X(08)71257-6}{DOI: 10.1016/S0049-237X(08)71257-6}. Principles for mechanisms from which computability follows.
\par\medskip

\noindent Douglas Harper (n.d.). Online Etymology Dictionary, entry ``ledger''. \href{https://www.etymonline.com/word/ledger}{Entry}. A large book that lay in one place, around 1400; the account-book sense from the 1580s.
\par\medskip

\noindent Clay Mathematics Institute (n.d.). The Millennium Prize Problems. \href{https://www.claymath.org/millennium-problems/}{Web page}. Announced at a meeting in Paris on 24 May 2000, a million dollars each; P versus NP is listed as unsolved and Navier--Stokes as active.
\par\medskip

\noindent Clay Mathematics Institute (n.d.). Navier--Stokes Equation. \href{https://www.claymath.org/millennium/navier-stokes-equation/}{Web page}. The problem's status, listed as active.
\par\medskip

\noindent Clay Mathematics Institute (n.d.). The Poincar\'e Conjecture. \href{https://www.claymath.org/millennium/poincare-conjecture/}{Web page}. Perelman's preprints of 2002 and 2003, and the prize of 2010.
\par\medskip

\noindent Konstantin Kakaes (2026). AI Has Solved One of Math's \$1 Million Millennium Prize Problems. Quanta Magazine, 8 September 2026. \href{https://www.quantamagazine.org/ai-has-solved-one-of-maths-1-million-millennium-prize-problems-20260908/}{Article}. About ten thousand agents, eighty-eight hours, and a Lean formalization that took about seventeen hours more.
\par\medskip

\noindent Ming Li and Paul Vit\'anyi (1996). Reversibility and adiabatic computation: trading time and space for energy. Proceedings of the Royal Society of London A 452, pp. 769--789. \href{https://doi.org/10.1098/rspa.1996.0039}{DOI: 10.1098/rspa.1996.0039}; \href{https://arxiv.org/abs/quant-ph/9703022}{arXiv: quant-ph/9703022}. One unit of price for each bit erased or taken in irreversibly, and nothing for reversible steps.
\par\medskip

\noindent OpenAI (2026). Finite Time Blowup for Navier--Stokes. Preprint, 166 pages. \href{https://cdn.openai.com/pdf/32d9f210-8b73-45e0-91bc-82a30aef8a9a/navier-stokes.pdf}{PDF}. A smooth force, compactly supported, that drives a three-dimensional flow starting from rest to unbounded velocity in finite time; the paper says this establishes Fefferman's alternatives (C) and (D).
\par\medskip

\noindent OpenAI (2026). On the Navier--Stokes Millennium Prize Problem. Announcement, 8 September 2026. \href{https://openai.com/index/navier-stokes-solution/}{Web page}.
\par\medskip

\noindent Grisha Perelman (2002). The entropy formula for the Ricci flow and its geometric applications. Preprint. \href{https://arxiv.org/abs/math/0211159}{arXiv: math/0211159}.
\par\medskip

\noindent Grisha Perelman (2003). Ricci flow with surgery on three-manifolds. Preprint. \href{https://arxiv.org/abs/math/0303109}{arXiv: math/0303109}.
\par\medskip

\noindent Grisha Perelman (2003). Finite extinction time for the solutions to the Ricci flow on certain three-manifolds. Preprint. \href{https://arxiv.org/abs/math/0307245}{arXiv: math/0307245}.
\par\medskip

\noindent Andrew Chi-Chih Yao (1981). Should tables be sorted? Journal of the ACM 28(3), pp. 615--628. \href{https://doi.org/10.1145/322261.322274}{DOI: 10.1145/322261.322274}. The cell-probe model, which counts only the looks into memory; an earlier version appeared at FOCS 1978.
\par\medskip

\subsection*{Afterword}

\noindent Meta (2021). Introducing Meta: A Social Technology Company. Meta Newsroom, 28 October 2021. \href{https://about.fb.com/news/2021/10/facebook-company-is-now-meta/}{Announcement}. Facebook's parent company takes the name Meta.
\par\medskip

\subsection*{What's mine, and what isn't}

\noindent Robert K. Merton (1957). Priorities in Scientific Discovery: A Chapter in the Sociology of Science. American Sociological Review 22(6), pp. 635--659. \href{https://doi.org/10.2307/2089193}{DOI: 10.2307/2089193}. Fights over priority in science.
\par\medskip

\noindent Stephen M. Stigler (1980). Stigler's Law of Eponymy. Transactions of the New York Academy of Sciences 39(1 Series II), pp. 147--157. \href{https://doi.org/10.1111/j.2164-0947.1980.tb02775.x}{DOI: 10.1111/j.2164-0947.1980.tb02775.x}. No discovery is named after the person who first found it, credited to Merton.
\par\medskip

